\documentclass[11pt,a4paper]{article}
\usepackage[utf8]{inputenc}
\usepackage[T1]{fontenc}
\usepackage{amsmath,amssymb,amsthm}
\usepackage{booktabs,array,caption,placeins,microtype,xcolor}
\usepackage{pgfplots}\pgfplotsset{compat=1.18}
\usepackage[margin=2.5cm]{geometry}
\usepackage[colorlinks=true,linkcolor=blue!60!black,citecolor=blue!60!black,urlcolor=blue!60!black]{hyperref}
\newtheorem{theorem}{Theorem}
\newtheorem{proposition}[theorem]{Proposition}
\newtheorem{lemma}[theorem]{Lemma}
\newtheorem{conjecture}[theorem]{Conjecture}
\newtheorem{corollary}[theorem]{Corollary}
\newtheorem{hypothesis}[theorem]{Hypothesis}
\theoremstyle{remark}\newtheorem{remark}[theorem]{Remark}
\newcommand{\Z}{\mathbb Z}\newcommand{\R}{\mathbb R}\newcommand{\Q}{\mathbb Q}\newcommand{\C}{\mathbb C}
\newcommand{\e}{\mathrm e}\newcommand{\Off}{\mathrm{Off}}\newcommand{\rep}[2]{\langle #1\rangle_{#2}}
\newcommand{\todo}[1]{\textcolor{red!70!black}{\small\sffamily[#1]}}

\title{The pair singular series of a polynomial IV:\\ the second spectrum}
\author{Jasper Reichardt\thanks{Independent researcher. Correspondence: jasperreichardt@gmail.com. This work was
carried out with heavy use of a large language model; see the disclosure at the end.}}
\date{Working draft of 13 September 2026 (corrected)}

\begin{document}
\maketitle

\begin{abstract}
Part~III reduced the Ces\`aro form of the off-diagonal of the pair singular series of an irreducible quadratic $f$
to ``pieces'' $\mathcal P_u(Y)$, one for each way the two roots of $f$ can agree on a part $u$ of the modulus, and
showed that everything except the pieces' ``windows'' is understood. This paper is about what the pieces actually
are. Numerically each piece is $O(\sqrt Y)$, and $\mathcal P_u(Y)/\sqrt Y$ oscillates in $\log Y$ at the Laplace
eigenvalues of the \emph{even} Maass cusp forms of $\mathrm{SL}_2(\Z)$, with amplitudes governed by the values of
the forms at the Heegner points, or their integrals over the closed geodesics, of discriminant $D$: a second
spectrum, next to the zeros of $\zeta_K$ that govern the diagonal (part~II). We explain this by an explicit
formula: the piece is a twisted Walfisz sum whose Dirichlet series is a sum over frequencies of Dirichlet series
of Sali\'e sums, which the Kuznetsov formula of half-integral weight continues to the half-line with poles at
$\tfrac12\pm it_j$, and the Katok--Sarnak correspondence identifies the residues as periods of Maass forms. The same law, and the same parameter-free
prediction of its phases, holds for a barer and more classical object: the partial sums of the Weyl sums of the
congruence $b^2\equiv D$, which Hooley bounded by $X^{3/4+\varepsilon}$ and which we find to exhibit square-root
cancellation with an explicit spectral expansion. We prove it for $u=1$ on a model object: the root pairs are Heegner points, the piece is a
sum of one test function over a modular orbit, the convergence obstruction and the passage to $L^2$ are resolved by a
truncation and by the symmetry of the orbit, the constant and continuous contributions vanish or are negligible
because the seed has vanishing horizontal mean, the arising Bessel integral is evaluated exactly by a Mellin--Barnes
contour shift, and the spectral sum converges absolutely for Riesz means of order at least two. The argument runs at
the levels $\Gamma_0(e)$, $e\mid\operatorname{rad}(2D)$: the coprimality condition $(d,2D)=1$ carried by the moduli is not
invariant under $\mathrm{SL}_2(\Z)$, an error of the first version of this paper found by an external reader, but its
M\"obius resolution into the sub-families $\{e\mid d\}$ is invariant under $\Gamma_0(e)$, and the spectrum accordingly
contains the even Maass newforms of those levels next to the level-one forms, as the data confirm (for $D=-3$ the
largest line is the first even newform of level $3$). We discuss the uniformity in $u$ that the Ces\`aro
conjecture of part~I would need, and report the numerical evidence, including phase predictions with no free
parameter that agree to a median of $0.012$ radians.
\end{abstract}

\medskip\noindent\fbox{\parbox{\dimexpr\textwidth-2\fboxsep-2\fboxrule\relax}{\small\textbf{Correction notice (13 September 2026).}
An external assessment of the repository (revision a372a63, 12 September 2026; file \texttt{research/reviews/}) found
that the proof of Theorem~\ref{thm:main} in the version of 12 September replaced the restricted sum defining the model
object~\eqref{eq:model}, whose moduli satisfy $(d,2D)=1$, by a sum over complete $\mathrm{SL}_2(\Z)$-orbits of Heegner
points, and that the restriction is not orbit-invariant: for $D=-4$ the admitted form $[5,4,1]$ and the excluded form
$[2,2,1]$ lie in the same orbit. We verified this; the theorem was not established as then stated. The present version
repairs the argument by working at the levels $\Gamma_0(e)$, $e\mid\operatorname{rad}(2D)$ (Lemma~\ref{lem:subfamily}
and Step~2 of \S\ref{sec:formula}), which also required moving the parametrisation from discriminant $D$ to
discriminant $4D$ (Step~1); every analytic step is unchanged, the level being fixed by $D$. The repaired theorem predicts
the even Maass newforms of the levels $e$ in the spectrum of the model object and their absence from the unrestricted
divisor sum, and this was checked numerically before the repair was written (Remark~\ref{rem:coprime}). The repair has
not been read by anyone outside; Remark~\ref{rem:referee} says what to check first. Record: \texttt{research/ERRATA.md}, item~13.}}

\section{Introduction}\label{sec:intro}

\paragraph{In plain terms.} Parts~I and~II showed that the main part of the fluctuation of prime values of a
polynomial is governed by the zeros of a Dedekind zeta function. The remaining part, the ``off-diagonal'' of
part~III, turns out to be governed by a second, entirely different spectrum: the Laplace eigenvalues of the
modular surface, the same numbers that govern the distribution of the roots of $x^2\equiv D\pmod d$ (Hooley,
Bykovski\u\i, Duke--Friedlander--Iwaniec). Only half of them appear, those of the even forms, and the size of each
contribution is the value of the corresponding eigenfunction at a special point, or its integral along a special
closed curve, determined by the discriminant of the polynomial. We found this in the data before we had the
theory; the theory, once written down, predicts exactly the pattern of presence and absence that the data show.

\paragraph{Set-up.} Notation as in part~III: $f=t^2+bt+c$, $D=b^2-4c$, $\lambda(p)=p/(p-2\omega(p))$; for squarefree
$u$ with $\omega(u)\ge1$ and $Q_u(x)=u^2x^2-D$,
\[
\mathcal P_u(Y)=\sum_{h\le Y}(Y-h)\bigl(F_u(Q_u(h))-\mathbb EF_u\bigr)-\tfrac12(\mathbb EF_u-1)Y,\qquad
F_u(n)=\sum_{\substack{d\mid n,\ d>1\\ d\ \mathrm{admissible}}}\frac{\lambda(d)}d,\qquad \mathbb EF_u=\sum_{\substack{d>1\\ d\ \mathrm{adm}}}\frac{\lambda(d)\omega(d)}{d^2},
\]
admissible meaning squarefree, all primes split, coprime to $2Du$. Part~III, Theorem~A$'$:
$\Off^\ast_f(H)=c_{\mathrm{off}}(f)H+\sum_{u\le H^{2/3+\varepsilon}}w(u)\mathcal W_u(H/u;\log H)+O(H(\log H)^{1-c})$, and the
Ces\`aro form of Conjecture~1 of part~I is equivalent to the windows being $o(H\log H)$.

\paragraph{Results.}
\begin{itemize}
\item \emph{The spectral formula for a piece} (Theorem~\ref{thm:main}; repaired on 13 September 2026, see the correction
notice): for $u=1$, fixed $f$, Riesz means of order $m\ge2$, and the model object~\eqref{eq:model}
of \S\ref{sec:numerics} rather than the piece itself (the weight $\lambda$ and the squarefree condition are dropped
there, the condition $(d,2D)=1$ is kept),
\[
\mathcal P^{(m)}_u(Y)=c_{u}Y+Y^{1/2}\sum_{j}\bigl(\alpha_j(u,f)Y^{it_j}+\overline{\alpha_j(u,f)}Y^{-it_j}\bigr)\,\gamma_m(t_j)+O(Y^{1/2-\delta}),
\]
the sum over the even Maass cusp forms $u_j$ of level one and the even newforms of the levels $e\mid\operatorname{rad}(2D)$,
with Laplace eigenvalue $\tfrac14+t_j^2$ (and, for $u>1$, the levels dividing $4u^2$), with $\alpha_j(u,f)$ an explicit
signed combination of periods of $u_j$ over Heegner points of level $e$ and discriminant $4D$, the Katok--Sarnak
periods for level one, and $\gamma_m(t)\ll|t|^{-m-1}$.
\item \emph{Numerical confirmation} (\S\ref{sec:numerics}): for five quadratics the even parameters explain the
variance of $\mathcal P_u(Y)/\sqrt Y$ at the $92$nd--$100$th percentile among random frequency sets, the odd ones at the
$12$th--$66$th; the dominant line $t_1=13.7798$ has stable amplitude and phase over two decades of $\log Y$; the two
polynomials with exponentially small periods ($D=-8,-163$) show nothing.
\item \emph{What the Ces\`aro conjecture needs} (\S\ref{sec:uniformity}): the sum of the pieces over $u\le H^{2/3}$
converges to $c_{\mathrm{off}}H+o(H)$ provided the spectral formula holds with an error $O((H/u)^{1/2+\varepsilon}u^{A})$,
$A<\tfrac14$; the level dependence of the Kuznetsov formula is therefore the whole remaining problem.
\end{itemize}

\section{The piece as a twisted Walfisz sum}\label{sec:walfisz}

Throughout, $f$ is an irreducible quadratic of discriminant $D$ and $u$ is squarefree with $\omega(u)\ge1$. Call a
modulus $d$ \emph{admissible} (for $u$) if it is squarefree and coprime to $2Du$. Only such $d$ occur: if
$p\mid Q_u(h)=u^2h^2-D$ and $p\nmid2Du$ then $D\equiv(uh)^2\pmod p$, so $\bigl(\tfrac Dp\bigr)=1$ and $p$ is split.
The condition ``every prime dividing $d$ is split'', which appears in part~III, is therefore automatic here, and we do
not impose it. For admissible $d$ put
\[
R_d=R^{(u)}_d=\{x\bmod d:\ u^2x^2\equiv D\},\qquad \rho(d)=|R_d|=\omega(d)=2^{\nu(d)},\qquad
\rho_k(d)=\sum_{x\in R_d}\e\Bigl(\frac{kx}d\Bigr),
\]
$\nu$ the number of prime factors and $\e(y)=e^{2\pi iy}$. Two properties are used constantly: $R_d$ is
\emph{symmetric}, $x\in R_d\Rightarrow-x\in R_d$, and $0\notin R_d$, both because $(D,d)=1$; consequently every
$\rho_k(d)$ is real. Write $\lambda(p)=p/(p-4)$ on split primes, $\lambda$ multiplicative, and
\begin{equation}\label{eq:FE}
F_u(n)=\sum_{\substack{d\mid n,\ d>1\\ d\ \mathrm{admissible}}}\frac{\lambda(d)}d,\qquad
\mathbb EF_u=\sum_{\substack{d>1\\ d\ \mathrm{admissible}}}\frac{\lambda(d)\rho(d)}{d^2},\qquad Q_u(h)=u^2h^2-D,
\end{equation}
both sums absolutely convergent ($\lambda(d)\rho(d)\ll d^{\varepsilon}$). \emph{All sums over moduli below run over
admissible $d>1$}: including $d=1$ would add $1$ to $F_u$ and to $\mathbb EF_u$, leaving~\eqref{eq:Sm} unchanged but
spoiling the symmetry used in Lemma~\ref{lem:sawtooth}, since $R_1=\{0\}$ is not a union of pairs $\{x,-x\}$ with
$x\ne-x$. The \emph{sharp piece} and the
\emph{Riesz pieces} are
\begin{equation}\label{eq:Sm}
S_u(t)=\sum_{h\le t}\bigl(F_u(Q_u(h))-\mathbb EF_u\bigr),\qquad
\mathcal P^{(m)}_u(Y)=\frac1{m!}\sum_{h\le Y}(Y-h)^m\bigl(F_u(Q_u(h))-\mathbb EF_u\bigr)\quad(m\ge0),
\end{equation}
so that $\mathcal P^{(0)}_u=S_u$ and, for integers $Y$, $\mathcal P^{(1)}_u(Y)=\sum_{t<Y}S_u(t)$; up to the explicit
constant of part~III, $\mathcal P^{(1)}_u(Y)$ is the piece $\mathcal P_u(Y)$ of part~III.
The connection with part~III is that $F_u(Q_u(h))$ is the divisor-sum form of the inner sum of a piece: $d\mid Q_u(h)$
exactly when $h$ lies in one of the $\rho(d)$ classes $R_d$.

\begin{lemma}[sawtooth form]\label{lem:sawtooth}
Let $\psi(y)=\tfrac12-\{y\}$. For every $t\ge0$,
\begin{equation}\label{eq:sawtooth}
S_u(t)=\sum_{d\ \mathrm{admissible}}\frac{\lambda(d)}d\sum_{x\in R_d}\psi\Bigl(\frac{t-x}d\Bigr),
\end{equation}
the series converging absolutely.
\end{lemma}

\begin{proof}
The right-hand side is not absolutely convergent term by term, since $\sum_d\lambda(d)\rho(d)/d$ diverges; it is summed
with the terms $x$ and $-x$ grouped, and converges in that sense, as shown at the end. Expanding $F_u$ and counting,
\[
\sum_{h\le t}F_u(Q_u(h))=\sum_{d}\frac{\lambda(d)}d\,\#\{h\le t:\ d\mid Q_u(h)\}
=\sum_{d}\frac{\lambda(d)}d\sum_{x\in R_d}\#\{h\le t:\ h\equiv x\ (d)\},
\]
the interchange being legitimate because $\sum_{d\mid Q_u(h)}\lambda(d)/d$ converges absolutely for each $h\le t$ and
there are finitely many $h$. Represent each $x\in R_d$ by $\rep xd\in[1,d]$. For $0\le t$ and $1\le x\le d$,
\[
\#\{1\le h\le t:\ h\equiv x\ (d)\}=\Bigl\lfloor\frac{t-x}d\Bigr\rfloor+1,
\]
valid also when $x>t$, where both sides vanish because $-1<(t-x)/d<0$. Writing $\lfloor y\rfloor=y-\tfrac12+\psi(y)$,
\[
\#\{h\le t:\ h\equiv x\ (d)\}=\frac td+\Bigl(\frac12-\frac xd\Bigr)+\psi\Bigl(\frac{t-x}d\Bigr).
\]
Now sum over $x\in R_d$. Since $R_d$ is symmetric and $0\notin R_d$, its representatives split into pairs
$\{\rep xd,\,d-\rep xd\}$, and $(\tfrac12-\tfrac xd)+(\tfrac12-\tfrac{d-x}d)=0$; hence
$\sum_{x\in R_d}(\tfrac12-\rep xd/d)=0$ for every admissible $d>1$. Therefore
\[
\sum_{h\le t}F_u(Q_u(h))=t\sum_{d}\frac{\lambda(d)\rho(d)}{d^2}+\sum_d\frac{\lambda(d)}d\sum_{x\in R_d}\psi\Bigl(\frac{t-x}d\Bigr)
=t\,\mathbb EF_u+\sum_d\frac{\lambda(d)}d\sum_{x\in R_d}\psi\Bigl(\frac{t-x}d\Bigr),
\]
which is~\eqref{eq:sawtooth}. For the convergence, group the terms of a pair $\{x,-x\}$ and write $a=\rep xd$ with
$a\le d/2$, so that the pair is $\{a,d-a\}$. If $d>2t$ and $a>t$ then $\{(t-a)/d\}=(t-a+d)/d$ and
$\{(t+a)/d\}=(t+a)/d$, so by the periodicity of $\psi$ the two sawtooths sum to $-2t/d$; if $a\le t$ their sum is
bounded by $1$, and this happens for at most $\rho(d)$ values with $a\le t<d/2$, a set that is empty once $d>2t$.
Hence for $d>2t$ each pair contributes $\ll t/d$ and
\[
\sum_{d>2t}\frac{\lambda(d)}d\cdot\frac{t\,\rho(d)}{d}\ \ll\ t\sum_{d>2t}\frac{\lambda(d)\rho(d)}{d^2}\ \ll\ t^{\varepsilon},
\]
while the range $d\le2t$ is a finite sum. So the paired series converges and~\eqref{eq:sawtooth} holds.
\end{proof}

\emph{In words.} A piece counts, for each admissible modulus $d$, how far the number of $h\le t$ with $d\mid Q_u(h)$
departs from its expectation, weighted by $\lambda(d)/d$; and that departure is a sawtooth evaluated at the roots.
The linear term that normally accompanies such a count cancels because the roots come in pairs $\pm x$. This is the
same cancellation that makes the diagonal of part~I a sum of sawtooths $\{H/d\}-\tfrac12$, and it is why the piece is
a Walfisz-type sum: in the classical case $\sum_{n\le x}n^{-1}(\tfrac12-\{x/n\})$ one sums a sawtooth over all moduli
with weight $1/n$, and here one sums it over admissible moduli, at the roots of a quadratic congruence, with weight
$\lambda(d)/d$.

\begin{lemma}[Weyl-sum form]\label{lem:weyl}
For every $t\ge0$ and every $K\ge1$,
\begin{equation}\label{eq:weylform}
S_u(t)=\frac1\pi\sum_{k\ge1}\frac1k\sum_{d\ \mathrm{admissible}}\frac{\lambda(d)\rho_k(d)}{d}\,\sin\frac{2\pi kt}d ,
\end{equation}
the $k$-sum converging boundedly. Consequently, for the Riesz pieces of order $m\ge1$ and integer $Y$,
\begin{equation}\label{eq:rieszweyl}
\mathcal P^{(m)}_u(Y)=\frac1\pi\sum_{k\ge1}\frac1k\sum_{d}\frac{\lambda(d)\rho_k(d)}{d}\,\Phi_m\Bigl(\frac{2\pi k}d;Y\Bigr),
\qquad \Phi_m(\theta;Y)=\frac1{m!}\sum_{h\le Y}(Y-h)^m\sin(\theta h),
\end{equation}
with $\Phi_m(\theta;Y)\ll_m\min\bigl(Y^{m+1},\ Y^{m}\|\theta/2\pi\|^{-1}\bigr)$, $\|\cdot\|$ the distance to $\Z$.
\end{lemma}

\begin{proof}
The Fourier expansion $\psi(y)=\sum_{k\ge1}\sin(2\pi ky)/(\pi k)$ holds for $y\notin\Z$, boundedly and uniformly on
compact subsets of $\R\setminus\Z$, and the partial sums are uniformly bounded. Insert it into~\eqref{eq:sawtooth}
and use, for each $d$,
\[
\sum_{x\in R_d}\sin\frac{2\pi k(t-x)}d=\Im\Bigl[\e\Bigl(\frac{kt}d\Bigr)\overline{\rho_k(d)}\Bigr]=\rho_k(d)\,\sin\frac{2\pi kt}d,
\]
since $\rho_k(d)$ is real. The interchange of the $k$-sum with the $d$-sum is justified by the uniform boundedness of
the partial sums of the Fourier series together with the absolute convergence established in
Lemma~\ref{lem:sawtooth} (dominated convergence for series). For~\eqref{eq:rieszweyl}, apply $\frac1{m!}\sum_{h\le Y}(Y-h)^m$
to $t=h$ in~\eqref{eq:weylform}, which is a finite sum, and evaluate $\Phi_m$ by summation by parts $m$ times against
the geometric sum $\sum_{h\le Y}\e(\theta h/2\pi)\ll\min(Y,\|\theta/2\pi\|^{-1})$.
\end{proof}

\begin{remark}\label{rem:whyk}
Two features of~\eqref{eq:rieszweyl} decide everything that follows. First, the frequencies $k$ enter only through the
\emph{Weyl sums} $\rho_k(d)$ of the dilated roots, so the arithmetic of the problem is entirely contained in the
Dirichlet series
\begin{equation}\label{eq:Wk}
W_k(s)=\sum_{d\ \mathrm{admissible}}\frac{\lambda(d)\rho_k(d)}{d^{s}}\qquad(\Re s>1),
\end{equation}
one for each $k$. Second, the weight $1/k$ in front decays only like $k^{-1}$, so a bound for $W_k$ that loses a power
of $k$ is expensive: this is the same phenomenon that forced the condition $\theta+6B<1$ in part~III,
Theorem~A, and it is the reason we do not bound the $\rho_k(d)$ one frequency at a time. The route of this paper is
instead to continue $W_k(s)$ analytically, where the dependence on $k$ appears as a Fourier coefficient of an
automorphic form and therefore \emph{decays} on average.
\end{remark}

\begin{proposition}[the Mellin representation]\label{prop:mellin}
Let $m\ge1$ and let
\[
A_u(s)=\sum_{h\ge1}\frac{F_u(Q_u(h))-\mathbb EF_u}{h^{s}}\qquad(\Re s>1).
\]
Then $A_u$ extends to a function holomorphic in $\Re s>\tfrac12$ except possibly for the singularities inherited from
the $W_k$, and for integer $Y\ge2$ and any $c>1$,
\begin{equation}\label{eq:perron}
\mathcal P^{(m)}_u(Y)=\frac1{2\pi i}\int_{(c)}A_u(s)\,\frac{Y^{s+m}}{s(s+1)\cdots(s+m)}\,ds .
\end{equation}
Moreover, for $0<\Re s<1$,
\begin{equation}\label{eq:AviaW}
A_u(s)=\frac{2\Gamma(1-s)}{(2\pi)^{1-s}}\,\sin\frac{\pi s}2\;\sum_{k\ge1}k^{s-1}\,W_k(s+1)\;+\;\bigl(\text{entire}\bigr),
\end{equation}
whenever the $k$-sum converges, where $W_k$ is as in~\eqref{eq:Wk}.
\end{proposition}

\begin{proof}
\eqref{eq:perron} is the Riesz--Perron formula, valid because $F_u(Q_u(h))-\mathbb EF_u\ll h^{\varepsilon}$ so that $A_u$
converges absolutely for $\Re s>1$. For~\eqref{eq:AviaW}, expand $F_u$ and sum over each residue class:
\[
\sum_{h\ge1}\frac{F_u(Q_u(h))}{h^s}=\sum_d\frac{\lambda(d)}{d^{1+s}}\sum_{x\in R_d}\zeta\Bigl(s,\frac{\rep xd}d\Bigr),
\]
$\zeta(s,\alpha)$ the Hurwitz zeta function. Since $\zeta(s,\alpha)$ has residue $1$ at $s=1$ for every $\alpha$, the
pole of the $d$-th term is $\rho(d)\lambda(d)d^{-1-s}$, and summing over $d$ these residues total $\mathbb EF_u$, which
is exactly the pole of $\mathbb EF_u\zeta(s)$: each bracket
$\lambda(d)d^{-1-s}\sum_{x}\zeta(s,\rep xd/d)-\lambda(d)\rho(d)d^{-2}\zeta(s)$ is entire, which is the statement that
$A_u$ continues past $s=1$. Hurwitz's formula,
\[
\zeta(s,\alpha)=\frac{2\Gamma(1-s)}{(2\pi)^{1-s}}\Bigl[\sin\frac{\pi s}2\sum_{k\ge1}\frac{\cos2\pi k\alpha}{k^{1-s}}
+\cos\frac{\pi s}2\sum_{k\ge1}\frac{\sin2\pi k\alpha}{k^{1-s}}\Bigr],
\]
valid for $\Re s<0$ and by continuation for $\Re s<1$ with the series interpreted as the analytic continuations of the
periodic zeta functions, gives after summing over the symmetric set $R_d$
\[
\sum_{x\in R_d}\zeta\Bigl(s,\frac{\rep xd}d\Bigr)=\frac{2\Gamma(1-s)}{(2\pi)^{1-s}}\sin\frac{\pi s}2\sum_{k\ge1}\frac{\rho_k(d)}{k^{1-s}},
\]
the sine term vanishing identically because $\sum_{x\in R_d}\sin(2\pi kx/d)=0$ by symmetry. Multiplying by
$\lambda(d)d^{-1-s}$ and summing over $d$ gives~\eqref{eq:AviaW}.
\end{proof}

\emph{In words.} The piece is determined by one Dirichlet series for each frequency, namely the Dirichlet series of
the Weyl sums of the roots of $u^2x^2\equiv D$; the symmetry of the root set kills the odd half of Hurwitz's formula,
so only the even half survives. That the odd half disappears here is the first appearance of the parity phenomenon
that will select the even Maass forms in \S\ref{sec:ks}, and it has the same source: the root set of a quadratic
congruence is invariant under $x\mapsto-x$.

\section{Sali\'e sums and the Dirichlet series of the Weyl sums}\label{sec:salie}

By Remark~\ref{rem:whyk} everything now depends on the series $W_k(s)=\sum_d\lambda(d)\rho_k(d)d^{-s}$. This section
identifies its coefficients as Kloosterman sums of half-integral weight, which is what makes an analytic continuation
possible. Write, for $c\ge1$ odd with $(D,c)=1$ and $h\in\Z$,
\begin{equation}\label{eq:Wh}
\mathcal W_h(D;c)=\sum_{b\bmod c,\ b^2\equiv D}\e\Bigl(\frac{hb}c\Bigr),
\end{equation}
the Weyl sum of the roots of the \emph{fixed} quadratic congruence, and let
$K_\chi(m,n;p)=\sum_{x\bmod p}^{*}\bigl(\tfrac xp\bigr)\e\bigl(\tfrac{mx+n\bar x}p\bigr)$ be the Kloosterman sum
twisted by the Legendre symbol, $\varepsilon_p=1$ or $i$ according as $p\equiv1$ or $3\pmod4$.

\begin{lemma}[the dilation is a change of frequency]\label{lem:bridge}
For admissible $d$ and $k\ge1$, with $\bar u$ the inverse of $u$ modulo $d$,
\[
\rho_k(d)=\mathcal W_{k\bar u}(D;d).
\]
\end{lemma}
\begin{proof}
$x\mapsto b=ux$ is a bijection of $R^{(u)}_d=\{x:u^2x^2\equiv D\}$ onto $\{b:b^2\equiv D\}$ because $(u,d)=1$, and
$kx\equiv k\bar ub\pmod d$.
\end{proof}

So the dilation by $u$ does not change the quadratic; it moves the frequency from $k$ to $k\bar u$, a frequency that
depends on the modulus. This is the exact point at which the pieces with $u>1$ leave the range of the theorems of
Duke, Friedlander and Iwaniec, which are stated for frequencies fixed independently of the modulus.

\begin{lemma}[Sali\'e]\label{lem:salie}
Let $p$ be an odd prime with $\bigl(\tfrac Dp\bigr)=1$ and let $h\in\Z$ with $p\nmid h$. Then
\begin{equation}\label{eq:salie}
\mathcal W_{2h}(D;p)=\frac{1}{\varepsilon_p\sqrt p}\,K_\chi\bigl(D,h^2;p\bigr).
\end{equation}
For odd squarefree $c=c_1c_2$ with $(c_1,c_2)=1$ and $(D,c)=1$,
\begin{equation}\label{eq:multiplicative}
\mathcal W_h(D;c)=\mathcal W_{h\overline{c_2}}(D;c_1)\,\mathcal W_{h\overline{c_1}}(D;c_2),
\end{equation}
so~\eqref{eq:salie} determines $\mathcal W_h(D;c)$ for every odd squarefree $c$ coprime to $2D$.
\end{lemma}
\begin{proof}
\eqref{eq:multiplicative} is the Chinese remainder theorem: $b\mapsto(b\bmod c_1,b\bmod c_2)$ is a bijection of the
roots modulo $c$ onto the pairs of roots, and $b/c\equiv b_1\overline{c_2}/c_1+b_2\overline{c_1}/c_2\pmod1$. For
\eqref{eq:salie}, substitute $b=2\bar hy$ in~\eqref{eq:Wh}, which is a bijection of the roots of $b^2\equiv D$ onto the
roots of $y^2\equiv h^2D\bar4$ and turns $\e(2hb/p)$ into $\e(2y/p)$; hence
$\mathcal W_{2h}(D;p)=\sum_{y^2\equiv h^2D\bar4}\e(2y/p)$. Sali\'e's evaluation of the twisted Kloosterman sum,
$K_\chi(m,n;p)=\varepsilon_p\sqrt p\,\bigl(\tfrac mp\bigr)\sum_{y^2\equiv mn\,(p)}\e\bigl(\tfrac{2y}p\bigr)$
\cite[\S4.5]{Iwaniec1997}, applied with $m=D$, $n=h^2$ gives~\eqref{eq:salie}, the factor $\bigl(\tfrac Dp\bigr)$
being $1$ because $p$ is split. \footnote{Both identities were verified numerically for all split $p\le37$, all
$D\in\{-4,-3,5,8,12\}$ and $h\in\{1,2,3,5\}$, and~\eqref{eq:multiplicative} for composite $c$; the reference for
Sali\'e's evaluation should be pinned to a precise statement of the normalisation.}
\end{proof}

\begin{remark}[what Lemma~\ref{lem:salie} does and does not give]\label{rem:achieved}
Combining Lemmas~\ref{lem:bridge} and~\ref{lem:salie}, for even $k$ and admissible $d$ the coefficient of $W_k(s)$
is, up to the elementary factor $\varepsilon_d^{-1}d^{-1/2}$, a Kloosterman sum of half-integral weight with the two
arguments $D$ and $(k\bar u/2)^2$: the first fixed by the polynomial, the second carrying the frequency and the
dilation. This is the shape to which the Kuznetsov formula of weight one half applies, and it is the machinery behind
every known estimate for these Weyl sums \emph{on average over the modulus} \cite{Hooley1963,DFI1995,DFI2012}.

It is, however, \emph{not} the route to the oscillations observed in \S\ref{sec:numerics}, and it is worth saying
why, because the distinction is easy to miss. Selberg's Kloosterman zeta function is normalised as
$Z_{m,n}(\sigma)=\sum_cS(m,n;c)c^{-2\sigma}$, with poles at $\sigma=s_j$, $\lambda_j=s_j(1-s_j)$
\cite[Theorem~1]{GoldfeldSarnak1983}, and their Remark~3 covers exactly our case, weight $\tfrac12$ on
$\Gamma_0(4N)$. Tracing that normalisation through Lemma~\ref{lem:salie} and Proposition~\ref{prop:mellin} places the
singularities of $A_u$ at $s=-\tfrac12+2it_j$ and predicts an oscillation of $\mathcal P^{(1)}_u(Y)/\sqrt Y$ at
frequency $2t_j$. That prediction is false: fitted at the frequencies $2t_j$ the even set explains between $2.1\%$
and $2.2\%$ of the variance and sits at the $43$rd to $69$th percentile of random frequency sets drawn from the same
band, against the $99.7$th to $100$th percentile at the frequencies $t_j$ (\S\ref{sec:numerics};
\texttt{scripts/piece-freq.ts}). The reason for the mismatch is that $Z_{m,n}$ sums over the \emph{modulus}, whereas
$\mathcal P_u$ sums over the \emph{argument} $h$. The two are different analytic problems, and only the second
produces the observed law.
\end{remark}

\paragraph{The argument-side expansion.} The route that does fit is Bykovski\u\i's \cite{Bykovskii1984}: write the
sum over $h\le Y$ as a finite combination of Poincar\'e series for $\mathrm{SL}_2(\Z)$ and apply the spectral
decomposition in weight $0$ and level $1$. There the parametrisation $\lambda_j=s_j(1-s_j)$, $s_j=\tfrac12+it_j$,
produces terms $Y^{s_j}=\sqrt Y\,Y^{it_j}$ directly, exactly as the error term in the hyperbolic lattice-point
problem does. This accounts at once for all four features of the data: the frequencies are the $t_j$ themselves, the
scale is $\sqrt Y$, the spectrum is that of weight $0$ and of the levels $e\mid\operatorname{rad}(2D)$ with the
level-one lines dominant (the newforms of level $e$ were seen in the data only after the first version of this paper
had been corrected, Remark~\ref{rem:coprime}), and the forms that can appear are those pairing
non-trivially with the Heegner points or the closed geodesics, which by
Proposition~\ref{prop:parity} and Corollary~\ref{cor:parity-e} are the even ones. What we prove is the following; the proof occupies
\S\S\ref{sec:formula}--\ref{sec:proof}, and the precise form, with the order of the Riesz mean, is
Theorem~\ref{thm:main}. (Its first version, at level one only, was wrong; see the correction notice and
Remark~\ref{rem:coprime}.)

\begin{theorem}[the spectral formula for a piece]\label{thm:K}
Let $u=1$ and let $\mathcal G$ be the model object~\eqref{eq:model}, smoothed by a Riesz mean of sufficiently high
order. There are coefficients $\alpha_j(D)$, supported on the even Maass cusp forms of the groups $\Gamma_0(e)$,
$e\mid\operatorname{rad}(2D)$ --- the forms of level one and the newforms of the levels $e$ --- with
$\lambda_j=\tfrac14+t_j^2$ and proportional to periods of $u_j$ over Heegner points of level $e$ and discriminant
$4D$ (the Katok--Sarnak periods when $e=1$), such that
\[
\mathcal G(Y)=\sqrt Y\sum_{j}\bigl(\alpha_j(D)Y^{it_j}+\overline{\alpha_j(D)}Y^{-it_j}\bigr)+\mathcal E(Y)+O\bigl(Y^{1/2-\delta}\bigr)
\]
for some $\delta>0$, the sum converging absolutely and $\mathcal E$ being the contribution of the continuous
spectrum, which is $o(\sqrt Y)$.
\end{theorem}

For fixed $D$ this is in substance the content of Bykovski\u\i's spectral treatment of the divisor sums
$\sum_{|n|\le P}\sigma_{-s}(n^2+D)$ \cite{Bykovskii1984}, which is where an asymptotic of this shape was first
obtained; what we require beyond it is the identification of the oscillating terms, their amplitudes and the error
exponent. The proof given in \S\ref{sec:proof} does not go through the Weyl-sum series $W_k$ at all; it works with
the test function directly, which is what makes it elementary. Lemmas~\ref{lem:bridge} and~\ref{lem:salie} are retained because the identification of the Weyl sums as
Sali\'e sums is what any \emph{quantitative} treatment of the remainder, or any attempt at the uniformity in $u$ of
\S\ref{sec:uniformity}, will have to use.

\section{Katok--Sarnak and the parity rule}\label{sec:ks}

The route of \S\ref{sec:salie} would place the poles of $W_k$ at $\tfrac12\pm it_j$, $t_j$ running over the weight-$\tfrac12$ spectrum, with
residues proportional to $\overline{\rho_j(|D|)}\rho_j((k/2)^2)$. Two classical inputs turn this into a statement
about the Maass forms of $\mathrm{SL}_2(\Z)$.

\paragraph{Shimura and Katok--Sarnak.} The Shimura correspondence sends the weight-$\tfrac12$ forms of level $4$ in
Kohnen's plus space to weight-$0$ Maass forms of level $1$ with the same spectral parameter, and the image consists of
the \emph{even} forms; the odd Maass forms correspond to weight $\tfrac32$. Katok and Sarnak
\cite{KatokSarnak1993} identify the Fourier coefficients of the weight-$\tfrac12$ form at a fundamental discriminant
$D$ with a period of the corresponding Maass form $u_j$: for $D<0$ the sum of $u_j$ over the Heegner points of
discriminant $D$, and for $D>0$ the integral of $u_j$ over the closed geodesics of discriminant $D$. Writing
$\mathrm{Per}_D(u_j)$ for that period,
\begin{equation}\label{eq:ks}
\rho_j(|D|)=\kappa(t_j,D)\,\mathrm{Per}_D(u_j)
\end{equation}
with an explicit factor $\kappa$. \footnote{Only $|\kappa|$, and only its dependence on $D$, is used below; the normalisation of $\mathrm{Per}$ should be pinned to and the normalisation of $\mathrm{Per}$ to
\cite{KatokSarnak1993}; only $|\kappa|$, and only its dependence on $D$, is used below.}

\begin{proposition}[parity]\label{prop:parity}
If $u_j$ is an odd Maass form then $\mathrm{Per}_D(u_j)=0$, for every discriminant $D$. Consequently only the even
Maass forms of $\mathrm{SL}_2(\Z)$ can contribute to the pieces.
\end{proposition}
\begin{proof}
Let $\iota(z)=-\bar z$, the reflection in the imaginary axis; $u_j$ is odd exactly when $u_j\circ\iota=-u_j$. For
$D<0$ the Heegner points of discriminant $D$ are the roots $z_Q\in\mathbb H$ of the integral binary quadratic forms
$Q=[a,b,c]$ of discriminant $D$, one for each class, and $\iota(z_{[a,b,c]})=z_{[a,-b,c]}$; since $[a,b,c]\mapsto[a,-b,c]$
is the inversion of the class group, it permutes the classes, so $\iota$ permutes the set of Heegner points. Hence
$\mathrm{Per}_D(u_j)=\sum_Qu_j(z_Q)=\sum_Qu_j(\iota z_Q)=-\mathrm{Per}_D(u_j)$, which forces
$\mathrm{Per}_D(u_j)=0$. For $D>0$ the same involution permutes the closed geodesics $S_Q$ attached to the classes,
and reverses none of the identifications used in defining the cycle integral, so the same computation applies to
$\mathrm{Per}_D(u_j)=\sum_Q\int_{\Gamma_Q\backslash S_Q}u_j\,ds$.
\end{proof}

\emph{In words.} The Heegner points of a given discriminant, and likewise the closed geodesics, are symmetric about
the imaginary axis; an odd Maass form is antisymmetric about it; so its period vanishes. This is the third appearance
of one and the same symmetry. In \S\ref{sec:walfisz} it made the linear term of the sawtooth identity vanish; in
Proposition~\ref{prop:mellin} it killed the odd half of Hurwitz's formula; here it removes the odd half of the
spectrum. All three are the statement that the roots of a quadratic congruence come in pairs $\pm x$.

\begin{corollary}[the shape of a piece]\label{cor:shape}
In Theorem~\ref{thm:K} the spectral sum runs over even Maass cusp forms only, those of $\mathrm{SL}_2(\Z)$ and the
newforms of the levels $e\mid\operatorname{rad}(2D)$. The coefficient of a line $t_j$ is a signed combination, over $e$,
of the periods of the corresponding forms over the sub-families $\mathcal H_e$ of Lemma~\ref{lem:subfamily} (Step~3 of
\S\ref{sec:formula}); for the unrestricted divisor sum, which is the term $e=1$ alone, it is proportional to the
Katok--Sarnak periods $\sum_{g^2\mid4D}\mathrm{Per}_{4D/g^2}(u_j)$.
\end{corollary}

\paragraph{The size of the periods, and a visibility threshold.} For $D<0$ of class number one the period is the single
value $u_1(z_D)$ at $z_D=(-b+i\sqrt{|D|})/2$. Through the Fourier expansion
$u_j(z)=2\sqrt y\sum_{n\ge1}a_j(n)K_{it_j}(2\pi ny)\cos(2\pi nx)$ this is governed by $K_{it_j}(2\pi y)$ with
$y=\sqrt{|D|}/2$. The Bessel function $K_{it}(x)$ of imaginary order is of size $e^{-\pi t/2}$ and oscillatory while
$x<t$, and only begins to decay in $x$ once $x>t$. Hence the periods of a fixed form $u_j$ are all of comparable size
for
\begin{equation}\label{eq:threshold}
|D|<\Bigl(\frac{t_j}\pi\Bigr)^2,
\end{equation}
and collapse exponentially beyond it. For the first even form, $t_1=13.7798$, the threshold is $|D|\approx19.2$. This
is a falsifiable prediction: the first line should be visible for the small discriminants and invisible for
$|D|\ge43$, while the higher lines, having larger $t_j$, should persist further. Section~\ref{sec:numerics} reports
what is observed.

\section{The spectral formula for a piece}\label{sec:formula}

We now give the argument-side route for $u=1$ and the model object~\eqref{eq:model}. The moduli of the model carry
the condition $(d,2D)=1$, and the first version of this section (12 September 2026) treated the sum as a sum over
complete $\mathrm{SL}_2(\Z)$-orbits of Heegner points, which it is not (Remark~\ref{rem:coprime}). The argument below
works instead at the levels $e\mid\operatorname{rad}(2D)$: the condition is resolved by M\"obius inversion into
sub-families $\{e\mid d\}$, each a finite union of orbits of $\Gamma_0(e)$, and the spectral theory of
$\Gamma_0(e)\backslash\mathbb H$ replaces that of the modular surface. Nothing analytic changes; the level is fixed,
depending on $D$ alone. Throughout $D<0$; for $D>0$ every occurrence of a Heegner point is replaced by a closed
geodesic and the argument is parallel.

\paragraph{Step 1: the root pairs are Heegner points, of discriminant $4D$.} A pair $(d,x)$ with $d\ge1$ and
$x\bmod d$, $x^2\equiv D\pmod d$, determines the integral form $[d,2x,(x^2-D)/d]$ of discriminant $4D$, and every
form $[d,b,c]$ of discriminant $4D$ with $d\ge1$ arises exactly once in this way: $b^2-4dc=4D$ forces $b$ even, $b=2x$
with $x^2\equiv D\pmod d$, and $b\bmod2d$ corresponds to $x\bmod d$. This is Gauss's parametrisation in the form used
by Hooley~\cite[\S3]{Hooley1963} and studied in general by Welsh~\cite{Welsh2022}, and it holds for \emph{every} $d$,
even or odd, coprime to $D$ or not; the forms need not be primitive. Writing $\mathcal Q_{4D}$ for the set of all
integral forms $Q=[a,b,c]$ of discriminant $4D$ with $a\ge1$, primitive or not, and $z_Q=(-b+\sqrt{4D})/(2a)$ for the root
of $Q$ in $\mathbb H$,
\begin{equation}\label{eq:gauss}
\{(d,x):\ d\ge1,\ x\bmod d,\ x^2\equiv D\ (d)\}\ \longleftrightarrow\ \mathcal Q_{4D},\qquad
(d,x)\mapsto z_{(d,x)}=\frac{-x+\sqrt D}{d}.
\end{equation}
A form of content $g$ is $g$ times a primitive form of discriminant $4D/g^2$ with the same root, so the points
$z_{(d,x)}$ are the Heegner points of the discriminants $4D/g^2$, $g^2\mid4D$, and they fall into
$\sum_gh(4D/g^2)$ orbits of $\Gamma=\mathrm{PSL}_2(\Z)$. \footnote{Verified numerically: reducing the forms
$[d,2x,\ast]$ for all $d\le80$ gives exactly $\sum_gh(4D/g^2)$ classes, namely $2,2,2,3,4,4,6,6,6$ for
$D=-4,-3,-7,-8,-11,-15,-20,-24,-23$. The first version of this paper parametrised the pairs by forms of discriminant
$D$ with $b^2\equiv D\ (4d)$; that is a bijection onto the same pairs when $d$ is odd and coprime to $D$
(Lemma~\ref{lem:conventions}), but it is not available for the moduli the sub-families below contain, and it produces
different points of $\mathbb H$.} The two coordinates of a point carry exactly the two data of a root pair:
\begin{equation}\label{eq:coords}
\Im z_{(d,x)}=\frac{\sqrt{|D|}}{d},\qquad \Re z_{(d,x)}=-\frac xd.
\end{equation}

\begin{lemma}[the sub-families are $\Gamma_0(e)$-stable]\label{lem:subfamily}
For $e\ge1$ let $\mathcal H_e=\{[a,b,c]\in\mathcal Q_{4D}:\ e\mid a\}$, the forms whose modulus is divisible by $e$. Then
$\mathcal H_e$ is stable under $\Gamma_0(e)$ and under the reflection $\iota:[a,b,c]\mapsto[a,-b,c]$, and it is a finite
union of $\Gamma_0(e)$-orbits. For $e=1$ it is all of $\mathcal Q_{4D}$.
\end{lemma}
\begin{proof}
For $\gamma=\left(\begin{smallmatrix}\alpha&\beta\\ \gamma'&\delta\end{smallmatrix}\right)\in\Gamma_0(e)$ the
leading coefficient of $Q\circ\gamma$ is $Q(\alpha,\gamma')=a\alpha^2+b\alpha\gamma'+c\gamma'^2\equiv a\alpha^2\pmod e$,
since $e\mid\gamma'$, and $(\alpha,e)=1$ because $\alpha\delta-\beta\gamma'=1$; so $e\mid a$ is preserved. The
reflection fixes $a$. For the finiteness, $\mathcal Q_{4D}$ is a finite union of $\Gamma$-orbits, $\Gamma_0(e)$ has
finite index in $\Gamma$, so each $\Gamma$-orbit is a finite union of $\Gamma_0(e)$-orbits, and $\mathcal H_e$ is a
union of some of them. \footnote{Checked numerically: for fifteen pairs $(D,e)$ with $e\mid\operatorname{rad}(2D)$ the
number of $\Gamma_0(e)$-orbits met by the forms with $d\le60$, $120$ and $240$ is the same, e.g.\ $3$ for $(-4,2)$, $2$
for $(-3,3)$, $6$ for $(-20,5)$, $10$ for $(-20,10)$; and $3000$ random elements of $\Gamma_0(e)$ preserve $e\mid a$ in
every case. $\mathcal H_e$ is \emph{not} $\Gamma$-stable: for $D=-4$,
$\left(\begin{smallmatrix}1&0\\1&1\end{smallmatrix}\right)$ carries $[5,4,1]$ to $[2,2,1]$.}
\end{proof}

\begin{corollary}[parity at level $e$]\label{cor:parity-e}
Let $u$ be a Maass cusp form of $\Gamma_0(e)$ with $u\circ\iota=-u$, $\iota(z)=-\bar z$. Then
$\sum_{Q\in\Gamma_0(e)\backslash\mathcal H_e}u(z_Q)/|\Gamma_0(e)_{z_Q}|=0$, where $\Gamma_0(e)_{z_Q}$ is the
stabiliser of $z_Q$ in $\Gamma_0(e)$.
\end{corollary}
\begin{proof}
$\iota$ normalises $\Gamma_0(e)$ (conjugation by $\iota$ negates the off-diagonal entries), so it permutes the
$\Gamma_0(e)$-orbits in $\mathcal H_e$, which it preserves by Lemma~\ref{lem:subfamily}, and $\iota(z_Q)=z_{\iota Q}$
with stabilisers of equal order. The sum therefore equals its own negative. This is Proposition~\ref{prop:parity} at
level $e$.
\end{proof}

The moduli of~\eqref{eq:model} are those with $(d,2D)=1$, and
$\mathbf 1_{(d,2D)=1}=\sum_{e\mid\operatorname{rad}(2D)}\mu(e)\,\mathbf 1_{e\mid d}$. The sub-families $\mathcal H_e$ with
$e\mid\operatorname{rad}(2D)$ are the objects the argument is about; there are $2^{\omega(2D)}$ of them. For odd
$e\mid D$ the condition $e\mid a$ forces $e\mid b$ (as $x^2\equiv D\equiv0\pmod e$ with $e$ squarefree), so these are
Heegner points of level $e$ in the sense of Gross, Kohnen and Zagier, with root $b\equiv0\pmod{2e}$.

\paragraph{Step 2: the piece is a signed sum of sums over $\Gamma_0(e)$-orbits.} The sawtooth form of
Lemma~\ref{lem:sawtooth} applies to the model object (it used neither $\lambda$ nor squarefreeness; the modulus $d=1$,
included in $\sigma^*_{-1}$, contributes $\psi(t)-\tfrac12=0$ at integer $t$), so that
\[
S(t)=\sum_{h\le t}\bigl(\sigma^*_{-1}(h^2-D)-\mathbb E\sigma^*_{-1}\bigr)=\sum_{(d,2D)=1}\frac1d\sum_{x\in R_d}\psi\Bigl(\frac{t-x}d\Bigr),
\]
the series summed with $x$ and $-x$ paired. Its terms are indexed by the pairs $(d,x)$ with $(d,2D)=1$, and by
M\"obius inversion
\[
S(t)=\sum_{e\mid\operatorname{rad}(2D)}\mu(e)\,S_e(t),\qquad S_e(t)=\sum_{e\mid d}\frac1d\sum_{x\in R_d}\psi\Bigl(\frac{t-x}d\Bigr),
\]
each $S_e$ again summed in pairs and convergent by the tail estimate of Lemma~\ref{lem:sawtooth}: the only unpaired
roots are $x=0$, which occurs for the finitely many $d\mid D$, and $x=d/2$ for even $d$, which contributes
$\psi(t/d-\tfrac12)=-t/d$ for $d>2t$, like a pair. \footnote{The regrouping was checked in exact rational arithmetic
for $D=-4$, $t=37.5$ and $1000.5$, $d\le6000$.} Substituting~\eqref{eq:coords}, each $S_e$ is a sum over the
sub-family $\mathcal H_e$ taken modulo translation ($x\bmod d$ is $\Re z\bmod1$):
\begin{equation}\label{eq:orbit}
S_e(t)=\frac{1}{\sqrt{|D|}}\sum_{z\in\mathcal H_e/\Z}\Phi_t(z),\qquad
\Phi_t(z)=\Im z\cdot\psi\Bigl(\frac{t\,\Im z}{\sqrt{|D|}}+\Re z\Bigr),
\end{equation}
since $\frac1d\psi\bigl(\frac{t-x}d\bigr)=\frac{\Im z}{\sqrt{|D|}}\,\psi\bigl(\frac{t\Im z}{\sqrt{|D|}}+\Re z\bigr)$ at
$z=z_{(d,x)}$. By Lemma~\ref{lem:subfamily} each $\mathcal H_e$ is a finite union of $\Gamma_0(e)$-orbits. This is the
whole content of the change of route: the object is not a sum over moduli carrying Kloosterman sums, but a signed
combination of sums of one fixed test function over orbits of the groups $\Gamma_0(e)$, $e\mid\operatorname{rad}(2D)$,
in $\mathbb H$, with the length $t$ entering only as a dilation parameter of the test function.

\begin{remark}[the coprimality condition is not $\Gamma$-invariant: the error of the first version]\label{rem:coprime}
The first version of this section (12 September 2026) wrote $S(t)$ itself, with its condition $(d,2D)=1$, as a sum over
complete $\Gamma$-orbits of Heegner points of discriminant $D$, and applied the spectral theory of
$\Gamma\backslash\mathbb H$. That is wrong: the leading coefficient of a form is not a class invariant (its values over
an orbit are the numbers the form represents), and the set of forms with $(a,2D)=1$ is a proper non-empty subset of
every orbit. For $D=-4$, where $h(D)=1$, the admitted form $[5,4,1]$ is carried by $z\mapsto z/(z+1)$ to the excluded
form $[2,2,1]$; for the eleven discriminants $-3\ge D\ge-56$ every class contains forms of both kinds among $a\le300$.
The error was found by an external assessment of the repository (12 September 2026; \texttt{research/reviews/}), not
by us, and is recorded in the repository's errata (item~13). What \emph{is} invariant is the divisibility $e\mid a$,
under $\Gamma_0(e)$ (Lemma~\ref{lem:subfamily}); hence the M\"obius inversion of Step~2 and the passage to the groups
$\Gamma_0(e)$, which is the form of the argument now given.

Two consequences are visible in the data and were not predicted by the first version: the spectrum of the model
object contains, besides the even Maass forms of level one, the even \emph{newforms} of the levels
$e\mid\operatorname{rad}(2D)$; and the unrestricted divisor sum $\sum_{d\mid h^2-D}d^{-1}$, which is the term $e=1$
alone and a finite union of $\Gamma$-orbits, does not. Both were checked on 13 September 2026, before the repair was
written (\texttt{scripts/piece-divset.ts}, \texttt{piece-level2.ts}, $Y\le10^7$; Maass parameters from the LMFDB):
after the six level-one even lines are removed, the restricted object and its difference from the unrestricted sum
carry the eight smallest even newforms of level $e$ at the $90$th--$100$th percentile against $400$ random sets of
eight frequencies ($D=-4,-8$ with $e=2$; $D=-3$ with $e=3$), the unrestricted sum at the $16$th--$78$th; and for
$D=-3$ the largest single line of the restricted object is at $5.085$, where the first even Maass newform of level $3$
has $t=5.0987$, while at that frequency the unrestricted sum has $R^2=0.01$. Not seen: level-$2$ lines for $D=-3$,
although $2\mid\operatorname{rad}(2D)$; the corresponding periods may be small, and we record the absence. The same
non-invariance, at the level $4u^2$ that grows with $u$, is the obstruction of \S\ref{sec:uniformity}: the sub-family
there is $\{2u^2\mid b\}$ at discriminant $4u^2D$.
\end{remark}

\begin{remark}[the Weyl series is a Poincar\'e-type series]\label{rem:poincare}
Before going on, it is worth recording what~\eqref{eq:orbit} says about the Dirichlet series of the Weyl sums
themselves, since it is the reason the spectrum can enter at all. Directly from the coordinates~\eqref{eq:coords},
for $\Re s>1$,
\begin{equation}\label{eq:poincare}
\sum_{z\in\mathcal H_e/\Z}(\Im z)^{s}\,\e(m\,\Re z)
=|D|^{s/2}\sum_{e\mid d}\frac1{d^{s}}\sum_{x^2\equiv D\,(d)}\e\Bigl(\frac{-mx}{d}\Bigr):
\end{equation}
\emph{the Dirichlet series of the Weyl sums of a quadratic congruence is a Poincar\'e-type series summed over the
Heegner points of discriminant $4D$, modulo translation} --- over all of them for the unrestricted Weyl sums
($e=1$), and in the signed combination over $e$ for the Weyl sums over moduli coprime to $2D$. The proof
below does not use~\eqref{eq:poincare}; it works with the test function directly, which is what makes it elementary.
We record the identity because it is the structural statement behind everything, and because \S\ref{sec:hooley}
tests the same law on the bare sums it describes.
\end{remark}

\paragraph{Step 3: unfolding and the period.} One point of bookkeeping decides the shape of everything that follows.
The arithmetic parametrisation of~\eqref{eq:gauss} is by pairs $(d,x)$ with $x$ taken \emph{modulo} $d$, and
$x\mapsto x+d$ is $z\mapsto z-1$; so the sum in~\eqref{eq:orbit} runs over the points of $\mathcal H_e$ \emph{modulo
translation}, that is over $\Gamma_\infty\backslash\Gamma_0(e)z_Q$ and not over $\Gamma_0(e)z_Q$, where
$\Gamma_\infty=\{z\mapsto z+n\}$ is the stabiliser of the cusp $\infty$, the same subgroup for every $e$. This is not a
normalisation: a full orbit contains $z_Q+n$ for every $n\in\Z$, all of the same height, and
$\sum_{\gamma\in\Gamma}(\Im\gamma w)^2$ diverges for that reason, whereas
$\sum_{\gamma\in\Gamma_\infty\backslash\Gamma}(\Im\gamma w)^2=E(w,2)$ converges. Accordingly we set
\[
P_e[\Phi_t](w)=\sum_{\gamma\in\Gamma_\infty\backslash\Gamma_0(e)}\Phi_t(\gamma w),
\]
a Poincar\'e series for $\Gamma_0(e)$ in the classical sense, so that
\[
S_e(t)=\frac1{\sqrt{|D|}}\sum_{Q\in\Gamma_0(e)\backslash\mathcal H_e}\frac{P_e[\Phi_t](z_Q)}{|\Gamma_0(e)_{z_Q}|},
\]
the sum over the finitely many orbits of Lemma~\ref{lem:subfamily}, $\Gamma_0(e)_{z_Q}$ the stabiliser (of order $1$,
$2$ or $3$). The spectral decomposition of $P_e[\Phi_t]$ on $L^2(\Gamma_0(e)\backslash\mathbb H)$, together with the
classical unfolding $\langle P_e[\Phi_t],u\rangle=\int_0^\infty\!\!\int_0^1\Phi_t(z)\overline{u(z)}\,\frac{dx\,dy}{y^2}$
--- the $x$-integral over a single period, which is exactly the integral computed in Step~4; the cusp $\infty$ of
$\Gamma_0(e)$ has width $1$ --- gives
\begin{equation}\label{eq:spectral}
S_e(t)=\frac1{\sqrt{|D|}}\Bigl[(\text{constant and Eisenstein terms})+\sum_{j}\langle\Phi_t,u^{(e)}_j\rangle\,\mathrm{Per}_{D,e}\bigl(u^{(e)}_j\bigr)\Bigr],\qquad
\mathrm{Per}_{D,e}(u)=\sum_{Q\in\Gamma_0(e)\backslash\mathcal H_e}\frac{u(z_Q)}{|\Gamma_0(e)_{z_Q}|},
\end{equation}
the sum over an orthonormal basis $\{u^{(e)}_j\}$ of the Maass cusp forms of $\Gamma_0(e)$, which may be taken to
consist of Hecke eigenforms: the newforms of the levels dividing $e$ and the oldforms they generate, the level-one
forms among them. For $e=1$ the geometric factor is the Katok--Sarnak period of \S\ref{sec:ks}, summed over the
discriminants $4D/g^2$; for $e>1$ it is the same kind of sum over the Heegner points of level $e$. By
Corollary~\ref{cor:parity-e} it vanishes for every odd $u^{(e)}_j$: the parity rule is now visible as the statement
that $\mathcal H_e$ is symmetric about the imaginary axis while an odd form is antisymmetric.

\paragraph{Step 4: where the frequencies come from.} No continuation theorem is needed here. Expand $\psi$ in
Fourier series; since $\psi$ is odd,
\[
\psi(a+x)+\psi(a-x)=\frac2\pi\sum_{k\ge1}\frac{\sin(2\pi ka)\cos(2\pi kx)}{k},
\qquad a=\frac{t\,\Im z}{\sqrt{|D|}},
\]
so that $\Phi^+_t(z)=\frac y\pi\sum_k k^{-1}\sin(2\pi ka)\cos(2\pi kx)$. Pair this against the Fourier expansion at
the cusp $\infty$ of an even Maass cusp form of $\Gamma_0(e)$, $u_j(z)=2\sqrt y\sum_{n\ge1}a_j(n)K_{it_j}(2\pi ny)\cos(2\pi nx)$.
The $x$-integral is $\int_0^1\cos(2\pi kx)\cos(2\pi nx)\,dx=\tfrac12\delta_{n,k}$, so
\begin{equation}\label{eq:innerprod}
\langle\Phi^+_t,u_j\rangle=\sum_{k\ge1}\frac{a_j(k)}{\pi k}\,I_k,\qquad
I_k=\int_0^\infty y^{-1/2}\sin(\beta_ky)\,K_{it_j}(\alpha_ky)\,dy,\quad \alpha_k=2\pi k,\ \beta_k=\frac{2\pi kt}{\sqrt{|D|}} .
\end{equation}
Substituting $y=w/\beta_k$ gives $I_k=\beta_k^{-1/2}\int_0^\infty w^{-1/2}\sin(w)K_{it_j}(\alpha_kw/\beta_k)\,dw$,
in which $\alpha_k/\beta_k=\sqrt{|D|}/t$ is small. Two standard facts finish the computation: the small-argument
expansion
\begin{equation}\label{eq:Ksmall}
K_{i\tau}(x)=\tfrac12\Bigl[\Gamma(-i\tau)\bigl(\tfrac x2\bigr)^{i\tau}+\Gamma(i\tau)\bigl(\tfrac x2\bigr)^{-i\tau}\Bigr]+O(x^2),
\end{equation}
and $\int_0^\infty w^{\sigma-1}\sin w\,dw=\Gamma(\sigma)\sin(\pi\sigma/2)$ for $0<\Re\sigma<1$, applied at
$\sigma=\tfrac12\pm it_j$. Hence
\begin{equation}\label{eq:asymp}
I_k=\beta_k^{-1/2}\Bigl[\Gamma(-it_j)\Gamma\bigl(\tfrac12+it_j\bigr)\sin\tfrac{\pi(\frac12+it_j)}2\Bigl(\frac{\sqrt{|D|}}{2t}\Bigr)^{it_j}\cdot\tfrac12+\overline{(\ \cdot\ )}\Bigr]+O\bigl((\sqrt{|D|}/t)^{2}\bigr).
\end{equation}
Every $t$-dependence is now explicit: $\beta_k^{-1/2}\propto t^{-1/2}$ and the bracket carries $t^{\mp it_j}$. So
$\langle\Phi^+_t,u_j\rangle=t^{-1/2}\bigl(A_j(D)\,t^{-it_j}+\overline{A_j(D)}\,t^{it_j}\bigr)+O(t^{-5/2})$, and
by~\eqref{eq:spectral} the same holds for $S_e(t)$ with $A_j$ multiplied by $\mathrm{Per}_{D,e}(u_j)$, hence for
$S(t)$ after the signed sum over $e$. Summing over
$t<Y$ turns $t^{-1/2\mp it_j}$ into $Y^{1/2\mp it_j}$, which is the observed law.

Two features of the data are read off~\eqref{eq:asymp} directly. The frequencies are the $t_j$ themselves, not
$2t_j$, because~\eqref{eq:Ksmall} produces $x^{\pm it_j}$ and not $x^{\pm2it_j}$. And the phase carries the factor
$(\sqrt{|D|}/2)^{it_j}$, so that changing the discriminant shifts the phase by exactly
$\tfrac{t_j}2\log|D|$ plus the sign of $\mathrm{Per}_D(u_j)$: this is precisely the prediction~\eqref{eq:phasepred}
tested in Tables~\ref{tab:phases} and~\ref{tab:weylphase}, and it is what those tables confirm to a median of
$0.012$ radians. \footnote{Both standard facts were checked numerically; \eqref{eq:Ksmall} to a relative $10^{-6}$ at
$x\le0.1$ for $t=13.78,17.74$.}

\begin{remark}[what this route does \emph{not} need]\label{rem:notneeded}
The derivation above uses no continuation theorem for any zeta or Poincar\'e series: only the Fourier expansions of
$\psi$ and of $u_j$, the small-argument behaviour of the Bessel function, and one elementary integral. This matters,
because our two earlier attempts to route the argument through such a theorem both failed. The first, through
Selberg's Kloosterman zeta function, predicted the frequencies $2t_j$ and was refuted by the data
(Remark~\ref{rem:achieved}). The second, through the meromorphic continuation of the Poincar\'e series of
\cite{GoldfeldSarnak1983}, does not apply: their $P_m(z,s)$ carries the factor $\e(m\gamma z)$, whose imaginary part
contributes an exponential decay $e^{-2\pi m\Im\gamma z}$ and puts the series in $L^2$, whereas the series arising
in~\eqref{eq:poincare} carries only $\e(m\,\Re\gamma z)$ and is not the same object. The identity
\eqref{eq:poincare} remains true and remains the reason the spectrum enters; what is not available is a citable
continuation of the particular series it produces.
\end{remark}

The same computation bears on the size of these sums. Formally,~\eqref{eq:asymp} makes every spectral term of $S(t)$
of size $t^{-1/2}$, so the Riesz mean is $O(Y^{1/2+\varepsilon})$ as soon as the sum over $j$ converges: square-root
cancellation, against the $X^{3/4+\varepsilon}$ of Hooley~\cite{Hooley1963} and the $X^{2/3+\varepsilon}$ of
Bykovski\u\i~\cite{Bykovskii1984} for the sharp sums. That is exactly what \S\ref{sec:hooley} observes, with
$T_1(X)/\sqrt X$ below $0.56$, $0.53$ and $0.86$ for $D=-4,-3,-8$ over $X\le10^6$. We do not state it as a
proposition, because the convergence of the sum over $j$ is the unfinished step below and there is no shortcut round
it: an earlier version of this paper deduced it instead from a continuation theorem that, on inspection, applies to a
different series (Remark~\ref{rem:notneeded}).

Theorem~\ref{thm:K} is proved, in the precise form of Theorem~\ref{thm:main}, in \S\ref{sec:proof}.

\paragraph{Step 5: why the Poincar\'e series converges.} As it stands $|\Phi_t(z)|\le\tfrac12\Im z$ and
$\sum_{\gamma\in\Gamma_\infty\backslash\Gamma_0(e)}\Im(\gamma w)$ diverges, so $P_e[\Phi_t]$ is not absolutely convergent as
it stands. The divergence is exactly that of the Eisenstein series at $s=1$. It is removed by the symmetry of the sub-family $\mathcal H_e$,
which we may as well build into the seed.

\begin{lemma}[the symmetrised seed decays quadratically]\label{lem:seed}
Write $z=x+iy$, $a=a_t(y)=ty/\sqrt{|D|}$, and
$\Phi^+_t(z)=\tfrac12\bigl(\Phi_t(z)+\Phi_t(-\bar z)\bigr)=\tfrac y2\bigl(\psi(a+x)+\psi(a-x)\bigr)$. Then
\begin{enumerate}
\item[(i)] if $\|x\|>a$, where $\|\cdot\|$ is the distance to $\Z$ and $a<\tfrac12$, then
$\Phi^+_t(z)=-ay=-ty^2/\sqrt{|D|}$ exactly;
\item[(ii)] $|\Phi^+_t(z)|\le y/2$ always, and the exceptional set $\{x:\|x\|\le a\}$ has measure $2a$ in $x$
modulo $1$.
\end{enumerate}
Moreover $\Phi^+_t$ may replace $\Phi_t$ in~\eqref{eq:orbit} without changing the value, because $\mathcal H_e$ is
stable under $z\mapsto-\bar z$ (Lemma~\ref{lem:subfamily}): the pair $(d,x)$ maps to $(d,-x)$.
\end{lemma}
\begin{proof}
$\psi$ is odd and, on any interval free of integers, satisfies $\psi(\theta+a)=\psi(\theta)-a$. If $\|x\|>a$ then
neither $x$ nor $-x$ crosses an integer when shifted by $a$, so
$\psi(x+a)+\psi(-x+a)=\psi(x)+\psi(-x)-2a=-2a$, giving (i); (ii) is $|\psi|\le\tfrac12$ and the definition of the
exceptional set. The last statement is $z_{(d,-x)}=-\overline{z_{(d,x)}}$ for $D<0$.
\end{proof}

So the seed is $O(t\,y^2)$ off a set of $x$-measure $O(ty)$, on which it is $O(y)$; both contributions to
$\sum_{\gamma\in\Gamma_\infty\backslash\Gamma_0(e)}|\Phi^+_t(\gamma w)|$ are then dominated by
$t\sum_{\gamma\in\Gamma_\infty\backslash\Gamma}(\Im\gamma w)^2=tE(w,2)$, which converges (the cosets of $\Gamma_0(e)$
are among those of $\Gamma$). The
arithmetic form of this convergence is the tail estimate already proved in Lemma~\ref{lem:sawtooth}. This is the
fourth and last appearance of one symmetry: the root set of a quadratic congruence is stable under $x\mapsto-x$, and
that single fact makes the linear term of the sawtooth identity vanish (Lemma~\ref{lem:sawtooth}), kills the odd half
of Hurwitz's formula (Proposition~\ref{prop:mellin}), annihilates the periods of the odd Maass forms
(Proposition~\ref{prop:parity}), and now renders the Poincar\'e series convergent.

\subsection{Completing the decomposition}\label{sec:complete}

Three things were left open above: the two conventions for the modulus, the sense in which the spectral
decomposition of Step~3 is legitimate, and the order of Riesz mean at which the sum over the spectrum converges. We
settle all three.

\begin{lemma}[the two conventions]\label{lem:conventions}
Let $d$ be odd and coprime to $D$. Reduction modulo $d$ is a bijection
\[
\{b'\bmod2d:\ b'^2\equiv D\ (4d)\}\ \longrightarrow\ \{x\bmod d:\ x^2\equiv D\ (d)\},
\]
so that for such $d$ the pairs of~\eqref{eq:gauss} are also parametrised by the forms $[d,b',\ast]$ of discriminant
$D$, with $\e(2k\,b'/(2d))=\e(kx/d)$. For even $d$, or $(d,D)>1$, no such bijection exists in general: for
$D\equiv5\pmod8$ there is no form of discriminant $D$ with even leading coefficient at all. This is why the
sub-families of Lemma~\ref{lem:subfamily} are taken in discriminant $4D$.
\end{lemma}
\begin{proof}
Since $(d,2)=1$, the Chinese remainder theorem splits $b'^2\equiv D\ (4d)$ into $b'^2\equiv D\ (d)$ and
$b'^2\equiv D\ (4)$; the latter is solvable because $D\equiv0$ or $1\pmod4$ and determines $b'$ modulo $2$, which
together with $b'$ modulo $d$ determines $b'$ modulo $2d$. The phases agree because $b'\equiv x\pmod d$. For the last
claim, $b'^2-4dc=D$ with $d$ even forces $b'^2\equiv D\pmod8$, and $5$ is not a square modulo $8$. \footnote{The
bijection was verified for $D\in\{-11,-8,-7,-4,-3,5,8,12\}$ and all odd $d\le200$ coprime to $2D$.}
\end{proof}

\begin{lemma}[truncation, and why the decomposition is legitimate]\label{lem:trunc}
Every point of $\mathcal H_e$ satisfies $\Im z=\sqrt{|D|}/d\le\sqrt{|D|}$. Consequently, replacing
$\Phi^+_t$ by its truncation $\Phi^\flat_t=\Phi^+_t\cdot\mathbf 1_{\{\Im z\le\sqrt{|D|}\}}$ changes neither side
of~\eqref{eq:orbit}. The truncated seed satisfies
\begin{enumerate}
\item[(i)] $|\Phi^\flat_t(z)|\le\sqrt{|D|}/2$, and $\Phi^\flat_t(z)=0$ for $\Im z>\sqrt{|D|}$;
\item[(ii)] $\Phi^\flat_t(z)\ll t\,(\Im z)^2$ as $\Im z\to0$, in the sense of Lemma~\ref{lem:seed};
\item[(iii)] $\int_0^1\Phi^\flat_t(x+iy)\,dx=0$ for every $y$.
\end{enumerate}
Hence $P_e[\Phi^\flat_t]$ is bounded on $\Gamma_0(e)\backslash\mathbb H$, lies in $L^2$, and its spectral decomposition is
valid.
\end{lemma}
\begin{proof}
(i) and (ii) are Lemma~\ref{lem:seed} and the truncation. (iii) holds because $\psi$ has mean zero over a period, so $\int_0^1[\psi(a+x)+\psi(a-x)]dx=0$ for every $y$; the
truncation does not depend on $x$ and so preserves this. For the last statement, if $\Im w>\sqrt{|D|}$ then the
identity coset contributes nothing and every other coset $\gamma\in\Gamma_\infty\backslash\Gamma_0(e)$ has
$\Im\gamma w\le1/\Im w<1/\sqrt{|D|}$, so the sum is dominated by $t\,(E(w,2)-(\Im w)^2)$, which is bounded; on the
rest of a fundamental domain the sum is dominated by $tE(w,2)$, bounded there. So $P_e[\Phi^\flat_t]$ is bounded,
hence square-integrable on the finite-volume quotient $\Gamma_0(e)\backslash\mathbb H$.
\end{proof}

\begin{remark}[the continuous spectrum]\label{rem:eisenstein}
Property (iii) also disposes of the Eisenstein contribution. The group $\Gamma_0(e)$ has finitely many cusps
$\mathfrak a$ (one for $e=1$, $2^{\omega(e)}$ for squarefree $e$), with Eisenstein series $E_{\mathfrak a}(z,s)$ whose
Fourier expansion at the cusp $\infty$ is $\delta_{\mathfrak a\infty}y^s+\varphi_{\mathfrak a\infty}(s)y^{1-s}$ plus
non-constant modes $\varphi_{\mathfrak a\infty}(n,s)\sqrt y\,K_{s-1/2}(2\pi|n|y)\e(nx)$
\cite[Theorem~3.4]{Iwaniec2002}. Unfolding gives
$\langle P_e[\Phi^\flat_t],E_{\mathfrak a}(\cdot,\tfrac12+ir)\rangle=\int_0^\infty\!\!\int_0^1\Phi^\flat_t\,\overline{E_{\mathfrak a}(\cdot,\tfrac12+ir)}\,d\mu$,
and the two constant terms do not depend on $x$, so by (iii) they contribute nothing, for every cusp. Only the
non-constant modes remain, and against those the computation of Step~4 is the same one with $it_j$ replaced by $ir$
and $a_j(n)$ by the coefficients $\varphi_{\mathfrak a\infty}(n,\tfrac12+ir)$, which are of polynomial size in $|n|$
and $|r|$. The continuous spectrum therefore contributes $\sqrt Y\int c(r)Y^{ir}\,dr$ to the Riesz mean, which is
$o(\sqrt Y)$ by the Riemann--Lebesgue lemma as soon as $c\in L^1$, the Riesz factor $\asymp|r|^{-m-1}$ supplying the
integrability. This is why the discrete lines are what one sees.
\end{remark}

\begin{proposition}[which Riesz order]\label{prop:riesz}
Normalise $u_j$ in $L^2(\Gamma_0(e)\backslash\mathbb H)$. For a newform of level $e'\mid e$ write
$a_j(n)=\rho_j(1)\lambda_j(n)$ with $\lambda_j$ the Hecke eigenvalues; an oldform $u(e''z)$ has the coefficients of $u$
at the indices divisible by $e''$, and the statements below hold for it with the constants of $u$. Then
in~\eqref{eq:asymp}
\begin{equation}\label{eq:Ajsize}
|A_j(D)|\ \asymp_D\ t_j^{-1/2+o(1)} ,
\end{equation}
and the Riesz mean of order $m$ multiplies the $j$-th term by a factor $\asymp t_j^{-m-1}$. Consequently the
spectral sum of Theorem~\ref{thm:K} converges absolutely as soon as
\[
m\ \ge\ 1\quad\text{(using }\|u_j\|_\infty\ll t_j^{5/12+\varepsilon}\text{ on compact sets, Iwaniec--Sarnak)},
\qquad\text{or}\qquad m\ \ge\ 2\quad\text{(using only }\|u_j\|_\infty\ll t_j^{1/2+\varepsilon}\text{)} .
\]
In particular the Ces\`aro mean, $m=1$, suffices, which is the mean at which the numerics of \S\ref{sec:numerics}
were computed.
\end{proposition}
\begin{proof}
For~\eqref{eq:Ajsize}: the $k$-sum $\sum_ka_j(k)k^{-3/2}=\rho_j(1)\sum_k\lambda_j(k)k^{-3/2}$ converges absolutely
and is $\rho_j(1)t_j^{o(1)}$, since $\lambda_j(n)\ll n^{7/64+\varepsilon}$ at every level~\cite{KimSarnak2003} and
$\tfrac32-\tfrac7{64}>1$. The
$\Gamma$-factors of~\eqref{eq:asymp} have modulus
\[
|\Gamma(-it_j)|\,\bigl|\Gamma(\tfrac12+it_j)\bigr|\,\bigl|\sin\tfrac{\pi(\frac12+it_j)}2\bigr|
\asymp\sqrt{\frac{2\pi}{t_j}}e^{-\pi t_j/2}\cdot\sqrt{2\pi}e^{-\pi t_j/2}\cdot\frac{e^{\pi t_j/2}}2
\asymp\frac{1}{\sqrt{t_j}}\,e^{-\pi t_j/2},
\]
while $|\rho_j(1)|^2\asymp_e\cosh(\pi t_j)/L(1,\mathrm{sym}^2u_j)$ with $t_j^{-\varepsilon}\ll L(1,\mathrm{sym}^2u_j)\ll t_j^{\varepsilon}$
\cite{HoffsteinLockhart1994} gives $|\rho_j(1)|\asymp e^{\pi t_j/2}t_j^{o(1)}$, the implied constants depending on the
(fixed) level.
The two exponentials cancel and~\eqref{eq:Ajsize} follows. For the Riesz factor,
$\frac1{m!}\int_0^Y(Y-t)^mt^{-1/2-it_j}dt=\frac{Y^{m+1/2-it_j}}{m!}B(m+1,\tfrac12-it_j)$ and
$B(m+1,\sigma)=\Gamma(m+1)\Gamma(\sigma)/\Gamma(m+1+\sigma)\asymp|\sigma|^{-m-1}$. Finally
$|\mathrm{Per}_{D,e}(u_j)|\le N_{D,e}\|u_j\|_\infty$, with $N_{D,e}$ the number of orbits in Lemma~\ref{lem:subfamily}
and the supremum over the compact set containing the finitely many points, so the $j$-th amplitude is
$\ll t_j^{5/12+\varepsilon}\cdot t_j^{-1/2+o(1)}\cdot t_j^{-m-1}$, and by the Weyl law for $\Gamma_0(e)$,
$\#\{j:t_j\le T\}\sim\mathrm{vol}(\Gamma_0(e)\backslash\mathbb H)\,T^2/(4\pi)$, the sum over $j$ is dominated by $\int^\infty t^{5/12-1/2-m-1}\,t\,dt=\int^\infty t^{-m-1/12}\,dt$,
which converges for $m>\tfrac{11}{12}$. With the trivial sup-norm bound the exponent is $\int^\infty t^{-m}dt$ instead.
\end{proof}

\paragraph{The expansion~\eqref{eq:Ksmall}, uniformly in the order.} One point remains: \eqref{eq:Ksmall} is an
asymptotic as $x\to0$ \emph{for fixed order}, whereas we apply it at $x=\sqrt{|D|}\,w/t$ with $t_j$ ranging over the
whole spectrum. What is needed is a version uniform in $t_j$. It is elementary, and follows from the convergent
series for $K_\nu$ rather than from any uniform-asymptotic machinery.

\begin{lemma}[uniform small-argument expansion]\label{lem:Kuniform}
For $\tau\ge1$ and $x>0$,
\[
K_{i\tau}(x)=\tfrac12\Bigl[\Gamma(-i\tau)\bigl(\tfrac x2\bigr)^{i\tau}+\Gamma(i\tau)\bigl(\tfrac x2\bigr)^{-i\tau}\Bigr]+E(x,\tau),
\qquad |E(x,\tau)|\ \le\ |\Gamma(i\tau)|\,\bigl(e^{x^2/(4\tau)}-1\bigr).
\]
In particular $|E(x,\tau)|\le|\Gamma(i\tau)|\,x^2/(2\tau)$ whenever $x^2\le2\tau$.
\end{lemma}
\begin{proof}
From $K_\nu=\frac\pi{2\sin\nu\pi}(I_{-\nu}-I_\nu)$, the series
$I_\nu(x)=\sum_{n\ge0}(x/2)^{2n+\nu}/(n!\,\Gamma(n+\nu+1))$ and $\Gamma(z)\Gamma(1-z)=\pi/\sin\pi z$,
\[
K_{i\tau}(x)=\tfrac12\Gamma(i\tau)\Gamma(1-i\tau)\sum_{n\ge0}\frac{(x/2)^{2n}}{n!}
\Bigl[\frac{(x/2)^{-i\tau}}{\Gamma(n+1-i\tau)}-\frac{(x/2)^{i\tau}}{\Gamma(n+1+i\tau)}\Bigr],
\]
convergent for every $x$. The term $n=0$ is the stated main term, using $\Gamma(1-i\tau)=-i\tau\,\Gamma(-i\tau)$. For
$n\ge1$,
\[
\Bigl|\frac{\Gamma(1\mp i\tau)}{\Gamma(n+1\mp i\tau)}\Bigr|=\frac1{|(1\mp i\tau)(2\mp i\tau)\cdots(n\mp i\tau)|}\le\tau^{-n},
\]
since each factor has modulus at least $\tau$. As $|\Gamma(i\tau)|=|\Gamma(-i\tau)|$, the tail is at most
$|\Gamma(i\tau)|\sum_{n\ge1}(x^2/4\tau)^n/n!=|\Gamma(i\tau)|(e^{x^2/(4\tau)}-1)$.
\end{proof}

\footnote{The bound was checked numerically at $30$ pairs with $\tau\in\{2,5,13.78,40,100,300\}$ and
$x\in\{0.5,1,2,4,8\}$; it holds in every case, with ratio between $0.0002$ and $0.995$.}

Lemma~\ref{lem:Kuniform} is exactly what~\eqref{eq:asymp} needs. In $I_k$ the argument is $x=cw$ with
$c=\alpha_k/\beta_k=\sqrt{|D|}/t$, the same for every $k$; splitting the $w$-integral at $W=\sqrt{t_j}/c$ and applying
the lemma below $W$ gives an error
\[
\ll|\Gamma(it_j)|\,\frac{c^2}{t_j}\int_0^Ww^{3/2}\,dw\ \asymp\ |\Gamma(it_j)|\,c^{-1/2}\,t_j^{-3/4},
\]
against a main term of size $\asymp|\Gamma(it_j)|c^{-1/2}$: a saving of $t_j^{-3/4}$, uniform in $j$, which
Proposition~\ref{prop:riesz} has ample room to absorb. Note also that the stationary point of the combined phase,
where $t_j\,d(\log x)/dx$ meets $1/c$, sits at $x=\sqrt{|D|}$, a fixed point independent of $t_j$ and well inside
the range where the lemma applies.

\subsection{The remaining estimates, and the proof}\label{sec:proof}

Two estimates finish the argument: an exact treatment of the integral $I_k$, and a bound for the tail removed by the
truncation. Both are elementary. Throughout, the truncation height of Lemma~\ref{lem:trunc} is taken to be $T=t$
rather than $\sqrt{|D|}$; this is permitted, since any height above $\sqrt{|D|}$ leaves the orbit sum unchanged,
and it costs only $\|P_e[\Phi^\flat_t]\|_2\ll\sqrt t$.

\begin{lemma}[the integral, exactly]\label{lem:mellin}
Let $\alpha,\beta>0$, $\tau\ge1$, and $I(\alpha,\beta,\tau)=\int_0^\infty y^{-1/2}\sin(\beta y)K_{i\tau}(\alpha y)\,dy$.
Then for any $\varepsilon\in(0,\tfrac12)$,
\[
I=\beta^{-1/2}\,\Re\Bigl[\Gamma(-i\tau)\,\Gamma\bigl(\tfrac12+i\tau\bigr)\sin\tfrac{\pi(\frac12+i\tau)}2\,
\Bigl(\frac\alpha{2\beta}\Bigr)^{i\tau}\Bigr]\;+\;O_\varepsilon\Bigl(\beta^{-1/2}\,e^{-\pi\tau/2}\,\tau^{A}\Bigl(\frac\alpha\beta\Bigr)^{2-\varepsilon}\Bigr)
\]
for an absolute constant $A$, uniformly in $\alpha,\beta,\tau$.
\end{lemma}
\begin{proof}
The Mellin transforms $\int_0^\infty y^{s-1}\sin(\beta y)\,dy=\Gamma(s)\sin(\pi s/2)\beta^{-s}$, valid for
$0<\Re s<1$, and $\int_0^\infty y^{s-1}K_{i\tau}(\alpha y)\,dy=2^{s-2}\alpha^{-s}\Gamma(\tfrac{s+i\tau}2)\Gamma(\tfrac{s-i\tau}2)$,
valid for $\Re s>0$, give by Parseval's formula for the Mellin transform, for any $c\in(0,\tfrac12)$,
\begin{equation}\label{eq:MB}
I=\frac1{2\pi i}\int_{(c)}\Gamma(s)\sin\frac{\pi s}2\,\beta^{-s}\;2^{-\frac32-s}\,\alpha^{s-\frac12}\,
\Gamma\Bigl(\frac{\frac12-s+i\tau}2\Bigr)\Gamma\Bigl(\frac{\frac12-s-i\tau}2\Bigr)ds .
\end{equation}
In $\Re s>0$ the integrand has poles only where the last two $\Gamma$-factors do, namely at
$s=\tfrac12\pm i\tau+2n$, $n\ge0$; $\Gamma(s)$ contributes none there and $\sin(\pi s/2)$ contributes none at all.
Move the contour to $\Re s=\tfrac52-\varepsilon$, crossing exactly the two poles with $n=0$. Their residues are the
displayed main term, and on the new line $|\beta^{-s}\alpha^{s-1/2}|=\beta^{-1/2}(\alpha/\beta)^{2-\varepsilon}$, which
is the stated saving. That the shifted integral converges, and that its size is
$e^{-\pi\tau/2}\tau^{A}$ times the above, follows from Stirling's formula: on $\Re s=\tfrac52-\varepsilon$ the two
$\Gamma$-factors have argument of real part $-1+\tfrac\varepsilon2$ and imaginary parts $\tfrac12(\pm\tau-\Im s)$, so
their product decays like $e^{-\frac\pi4(|\tau-\Im s|+|\tau+\Im s|)}\le e^{-\pi\tau/2}$ times a power of
$1+|s|+\tau$, while $\Gamma(s)\sin(\pi s/2)$ grows at most like a power times $e^{-\pi|\Im s|/2}\cdot e^{\pi|\Im s|/2}$;
the resulting integral in $\Im s$ converges absolutely. \footnote{The main term was checked against numerical quadrature:
for $\tau=13.7798$ and $\alpha/\beta=0.1,0.04,0.02,0.01$ the relative error is $0.0715,0.0099,0.0018,0.00035$, i.e.
$\asymp(\alpha/\beta)^2$ with a bounded constant, confirming both the main term and the exponent.}
\end{proof}

\begin{lemma}[the truncated tail]\label{lem:tail}
For $\alpha,\beta>0$, $\tau\ge1$ and $\alpha T\ge1$,
\[
\Bigl|\int_T^\infty y^{-1/2}\sin(\beta y)K_{i\tau}(\alpha y)\,dy\Bigr|\ \ll\ \alpha^{-1/2}(\alpha T)^{-1/2}e^{-\alpha T},
\]
uniformly in $\beta$ and $\tau$.
\end{lemma}
\begin{proof}
From $K_\nu(x)=\int_0^\infty e^{-x\cosh u}\cosh(\nu u)\,du$ one gets, for $\nu=i\tau$,
$|K_{i\tau}(x)|\le\int_0^\infty e^{-x\cosh u}du=K_0(x)$. \footnote{Checked at $25$ pairs, $\tau\le120$, $x\le100$: the
inequality holds in every case.} Hence the integral is at most $\int_T^\infty y^{-1/2}K_0(\alpha y)\,dy$, and
substituting $x=\alpha y$ and using $K_0(x)\ll x^{-1/2}e^{-x}$ for $x\ge1$ gives the claim.
\end{proof}

\begin{theorem}[the spectral formula for the model object]\label{thm:main}
Let $D<0$ be a discriminant, $u=1$, and let $\mathcal G^{(m)}(Y)$ denote the Riesz mean of order $m$ of the model
object~\eqref{eq:model}. For each $e\mid\operatorname{rad}(2D)$ let $\{u^{(e)}_j\}$ be an orthonormal basis of the even
Maass cusp forms of $\Gamma_0(e)$, with Laplace eigenvalues $\tfrac14+t_{j,e}^2$. For $m\ge2$ there are coefficients
$\alpha_{j,e}(D)$, proportional to the periods $\mathrm{Per}_{D,e}(u^{(e)}_j)$ of~\eqref{eq:spectral}, such that
\[
\mathcal G^{(m)}(Y)=Y^{m}\sqrt Y\sum_{e\mid\operatorname{rad}(2D)}\mu(e)\sum_{j}\bigl(\alpha_{j,e}(D)\,Y^{it_{j,e}}+\overline{\alpha_{j,e}(D)}\,Y^{-it_{j,e}}\bigr)
\;+\;\mathcal E(Y)\;+\;O\bigl(Y^{m+\frac12-\delta}\bigr)
\]
for some $\delta>0$, each inner sum converging absolutely, where $\mathcal E$ is the contribution of the continuous
spectra and satisfies $\mathcal E(Y)=o(Y^{m+1/2})$. The spectral parameters that occur are those of the even Maass cusp
forms of level one and of the even newforms of the levels $e\mid\operatorname{rad}(2D)$, and no others.
\end{theorem}
\begin{proof}
By Step~2, $S(t)=\sum_e\mu(e)S_e(t)$ is a finite signed sum and the Riesz mean is linear, so it suffices to treat one
$S_e$, given by~\eqref{eq:orbit} as a sum over the finitely many $\Gamma_0(e)$-orbits of $\mathcal H_e$
(Lemma~\ref{lem:subfamily}). Apply the Riesz mean of order $m$ to~\eqref{eq:orbit} and truncate the seed at height
$T=t$ (Lemma~\ref{lem:trunc}), which changes nothing. By Lemma~\ref{lem:seed} the automorphised seed
$P_e[\Phi^\flat_t]$ converges absolutely and lies in $L^2(\Gamma_0(e)\backslash\mathbb H)$, so it has a spectral
decomposition; by Proposition~\ref{prop:riesz} the Riesz factor $\asymp t_j^{-m-1}$ makes the discrete part converge
absolutely for $m\ge2$ even with only the trivial sup-norm bound, so the expansion converges uniformly and may be
evaluated at the points $z_Q$ and summed over the orbits, which replaces $u^{(e)}_j$ by
$\mathrm{Per}_{D,e}(u^{(e)}_j)$ (Step~3); the odd forms drop out by Corollary~\ref{cor:parity-e}. The constant term
vanishes because $\int_0^1\Phi^\flat_t\,dx=0$ (Lemma~\ref{lem:trunc}(iii)), and the same identity confines the
Eisenstein contribution of every cusp to the non-constant modes, whence $\mathcal E(Y)=o(Y^{m+1/2})$ by
Remark~\ref{rem:eisenstein}. The last sentence of the statement holds because the basis may be chosen to consist of
newforms of levels $e'\mid e$ and their oldform images, whose spectral parameters are those of the newforms.

For the coefficients, \eqref{eq:innerprod} expresses $\langle\Phi^\flat_t,u_j\rangle$ as $\sum_k(a_j(k)/\pi k)$
times $I(\alpha_k,\beta_k,t_j)$ minus the truncated tail, with $\alpha_k=2\pi k$ and $\beta_k=2\pi kt/\sqrt{|D|}$.
The interchange of the $k$-sum with the $y$-integral is justified by
$\int_0^\infty y^{-1/2}K_0(\alpha_ky)\,dy\ll k^{-1/2}$ together with $\sum_k|a_j(k)|k^{-3/2}<\infty$, which holds
because $\lambda_j(n)\ll n^{7/64+\varepsilon}$. Lemma~\ref{lem:mellin} evaluates each $I(\alpha_k,\beta_k,t_j)$ with a
relative error $O(\tau^{A}(\alpha_k/\beta_k)^{2-\varepsilon})$; crucially $\alpha_k/\beta_k=\sqrt{|D|}/t$ does not
depend on $k$, so the relative error is uniform in $k$ and survives the $k$-sum unchanged. Lemma~\ref{lem:tail}
bounds the tail by $\ll e^{-2\pi kt}$, which after summation over $k$ and, using the Riesz factor, over $j$, is
$O(e^{-2\pi t})$. Collecting the main terms gives
$\langle\Phi^\flat_t,u_j\rangle=t^{-1/2}(A_j(D)t^{-it_j}+\overline{A_j(D)}t^{it_j})+O(t^{-5/2+\varepsilon}t_j^{A})$
with $|A_j(D)|\asymp_Dt_j^{-1/2+o(1)}$ by~\eqref{eq:Ajsize}; the factor
$(\alpha_k/2\beta_k)^{it_j}=(\sqrt{|D|}/2t)^{it_j}$ supplies the dependence on the discriminant. Summing $t^{-1/2\mp it_j}$
against the Riesz weight turns it into $Y^{m+1/2\mp it_j}$ times $B(m+1,\tfrac12\mp it_j)\asymp t_j^{-m-1}$, and the
error terms sum by Proposition~\ref{prop:riesz} with room to spare, since they carry two further powers of $t$.
\end{proof}

\begin{remark}[what a referee should check]\label{rem:referee}
The proof uses, as standard inputs: Gauss's parametrisation of the root pairs by binary quadratic forms of
discriminant $4D$ (\S\ref{sec:formula}, Step~1), the spectral decomposition of $L^2(\Gamma_0(e)\backslash\mathbb H)$
for the finitely many levels $e\mid\operatorname{rad}(2D)$ with its Eisenstein series at each cusp, the Weyl law and the
sup-norm bound for Maass forms at fixed level, the Kim--Sarnak bound $\lambda_j(n)\ll n^{7/64+\varepsilon}$ and the
Hoffstein--Lockhart normalisation of $\rho_j(1)$ for newforms of fixed level, and Parseval's formula for the Mellin
transform. Everything else is proved above. The step that was wrong in the first version, the passage from the
restricted sum to orbits (Remark~\ref{rem:coprime}), was found by an external reader; its replacement is
Lemma~\ref{lem:subfamily} together with Step~2, and it is the first thing to check. The two further places we would
look are the uniformity in
$\tau$ of the shifted contour in Lemma~\ref{lem:mellin}, where we have given the Stirling estimate in outline rather
than in full, and the passage from the spectral expansion of $P_e[\Phi^\flat_t]$ in $L^2$ to its pointwise evaluation at
the Heegner points, which we justify by the absolute convergence supplied by the Riesz factor and which is the reason
the theorem is stated for $m\ge2$ rather than $m\ge1$. This paper has twice contained an argument that appealed to a
classical theorem about a different object (Remarks~\ref{rem:achieved} and~\ref{rem:notneeded}); we would rather name
the weak points than let a reader hunt for them.
\end{remark}

\section{Numerical evidence}\label{sec:numerics}

All computations are exact rational or high-precision arithmetic; the scripts are in the repository
(\texttt{research/paper-IV/scripts/}). For each $f$ and $u$ we compute $\mathcal P^{(1)}_u(Y)$ for all $Y\le10^7$
(a sieve of $Q_u(h)$ for $h\le Y$, with $\mathbb EF_u$ evaluated to $22$ significant digits, since an error
$\delta$ in $\mathbb EF_u$ enters as $\delta Y^2/2$), sample it on a logarithmic grid of $4096$ points in
$\log Y\in[\log10^3,\log Y_{\max}]$, remove the mean and the linear trend in $\log Y$, and regress the result on
\[
\bigl\{\cos(t_j\log Y),\ \sin(t_j\log Y)\bigr\}_{j=1}^{6}
\]
for the six smallest even spectral parameters, for the six smallest odd ones, and for $300$ random sets of six
frequencies drawn uniformly from $[12,25]$. Reported is the fraction of variance explained, $R^2$, and where it falls
among the random sets.

\begin{table}[htbp]
\centering\small
\caption{Fit of $\mathcal P^{(1)}_u(Y)/\sqrt Y$ to the even and to the odd spectral parameters of $\mathrm{SL}_2(\Z)$,
against $300$ random frequency sets of the same size ($Y\le10^7$; $u=1$ throughout except where shown).}\label{tab:fits}
\begin{tabular}{llrrrrr}
\toprule
$f$ & $D$ & $R^2$ even & percentile & $R^2$ odd & percentile & verdict\\
\midrule
$t^2-2$      & $8$    & $0.217$ & $100$ & $0.036$ & $34$ & even\\
$t^2+1$      & $-4$   & $0.152$ & $100$ & $0.028$ & $20$ & even\\
$t^2+t+1$    & $-3$   & $0.139$ & $100$ & $0.037$ & $35$ & even\\
$t^2+t+2$    & $-7$   & $0.089$ & $100$ & $0.051$ & $76$ & even\\
$t^2-3$      & $12$   & $0.066$ & $99$  & $0.037$ & $66$ & even\\
$t^2+t-1$    & $5$    & $0.062$ & $92$  & $0.036$ & $32$ & even\\
$t^2+t+5$    & $-19$  & $0.047$ & $92$  & $0.044$ & $88$ & inconclusive\\
$t^2+2$      & $-8$   & $0.031$ & $47$  & $0.054$ & $86$ & none\\
$t^2+t+3$    & $-11$  & $0.009$ & $44$  & $0.023$ & $100$& none\\
\bottomrule
\end{tabular}
\end{table}

Table~\ref{tab:fits} is the evidence for Proposition~\ref{prop:parity}. Six of the nine polynomials place the even set
at the $92$nd percentile or above, four of them above every one of the $300$ random sets; the odd set never rises above
the $88$th percentile except once, at the $100$th for $t^2+t+3$, where the even set is at the $44$th and the fit is
at the level of the noise. For the dominant line $t_1=13.7798$ the fitted amplitude and phase are stable when the
range of $\log Y$ is cut in half: for $t^2-2$, amplitudes $0.019$ and $0.018$ with phases $0.89$ and $0.61$; for
$t^2+t+1$, $0.021$ and $0.017$ with phases $1.01$ and $1.03$; for $t^2+1$ at $u=2$, $0.042$ and $0.042$ with phases
$-0.30$ and $0.07$. A spurious fit does not reproduce its phase on an independent range.

\paragraph{The periods.} Using the Fourier coefficients of the first even form (LMFDB label \texttt{1.0.1.3.1},
$t_1=13.7797513$) we evaluated $u_1$ at the Heegner points of the nine imaginary discriminants of class number one
(\texttt{scripts/maass-period.py}).

\begin{table}[htbp]
\centering\small
\caption{The Katok--Sarnak period of the first even Maass form at the Heegner point of discriminant $D<0$, class
number one, relative to $D=-4$; and the fitted amplitude of the line $t_1$, relative to $D=-4$.}\label{tab:periods}
\begin{tabular}{lrrrrrrrrr}
\toprule
$D$ & $-3$ & $-4$ & $-7$ & $-8$ & $-11$ & $-19$ & $-43$ & $-67$ & $-163$\\
\midrule
$|u_1(z_D)|/|u_1(i)|$ & $0.649$ & $1$ & $0.542$ & $0.318$ & $0.302$ & $0.444$ & $6.6\cdot10^{-3}$ & $1.0\cdot10^{-4}$ & $2.2\cdot10^{-10}$\\
fitted amplitude ratio & $0.483$ & $1$ & $0.111$ & $0.281$ & $0.461$ & $1.358$ & --- & --- & ---\\
signal present & yes & yes & yes & no & no & marginal & no & no & no\\
\bottomrule
\end{tabular}
\end{table}

Table~\ref{tab:periods} confirms the mechanism at the level of orders of magnitude and no further. The threshold
\eqref{eq:threshold} sits at $|D|\approx19.2$: below it every period lies within a factor $3$ of every other, and
every one of the six discriminants with $|D|\le19$ has an amplitude within a factor $12$ of the others; above it the
periods fall by two, four and ten orders of magnitude at $D=-43,-67,-163$, and no signal is seen at any of them. In
particular the absence at $D=-163$, the polynomial $t^2+t+41$, is explained. What the table does \emph{not} confirm is
the amplitude law itself: the two rows disagree by as much as a factor $5$ at $D=-7$. Two unmodelled effects are of
that size, the local factors, since $\mathbb EF_u$ ranges over $[0.35,1.60]$ across these polynomials and enters
multiplicatively, and the uncertainty of the fitted amplitudes themselves, which for $D=-7$ is a factor $2$ between
the two halves of the range. We record the amplitude law as a prediction to be tested, not as a confirmed fact.

\paragraph{The weight, and a model piece.} The absences at $D=-8$ and $D=-11$, whose periods are of ordinary size,
are not explained by anything above, and they are not isolated: among the fourteen polynomials we tested, the six
showing no signal are exactly those for which the prime $3$ splits in $\Q(\sqrt D)$. That is a statement about the
weight, not about the roots. Indeed $\lambda(3)=3/(3-4)=-3$ is the only negative value taken by $\lambda$ and the
largest in modulus, and $|\lambda(3)\rho(3)/3|=2>1$, so at $\Re s=\tfrac12$ the local factor at $3$ of the Euler
product for $W_k$ has modulus up to $3\cdot2\cdot3^{-1/2}=3.46$: it is the one local factor that can dominate the
spectral term. Removing the weight confirms this. Define the \emph{model piece} by replacing $\lambda$ by $1$ in
\eqref{eq:FE}, so that $F$ becomes the restricted divisor function $\sum_{d\mid n,\ d\ \mathrm{adm}}1/d$ and
$\mathcal P^{(1)}$ is a Riesz mean of $\sigma_{-1}$ along $Q_u$. Repeating the regression:

\begin{table}[htbp]
\centering\small
\caption{Removing the arithmetic weight, at $Y\le10^7$, $u=1$. Fit of $\mathcal P^{(1)}/\sqrt Y$ to the six even
spectral parameters, for the weighted piece of part~III, for the same sum with $\lambda\equiv1$, and for the sum over
all divisors rather than the squarefree ones. The odd sets, omitted, run from the $4.7$th to the $28.7$th percentile
throughout.}\label{tab:model}
\begin{tabular}{lrrrrrr}
\toprule
& \multicolumn{2}{c}{weighted, squarefree} & \multicolumn{2}{c}{$\lambda\equiv1$, squarefree} & \multicolumn{2}{c}{$\lambda\equiv1$, all divisors}\\
\cmidrule(lr){2-3}\cmidrule(lr){4-5}\cmidrule(lr){6-7}
$D$ & $R^2$ even & pct & $R^2$ even & pct & $R^2$ even & pct\\
\midrule
$8$   & $0.217$ & $100$  & $0.319$ & $100$  & $0.420$ & $100$\\
$-4$  & $0.152$ & $100$  & $0.237$ & $100$  & $0.361$ & $100$\\
$-3$  & $0.139$ & $100$  & $0.139$ & $100$  & $0.208$ & $100$\\
$-8$  & $0.031$ & $47$   & $0.109$ & $100$  & $0.162$ & $100$\\
$-11$ & $0.009$ & $44$   & $0.078$ & $100$  & $0.167$ & $100$\\
$17$  & $0.078$ & $99.3$ & $0.098$ & $99.3$ & $0.116$ & $100$\\
\bottomrule
\end{tabular}
\end{table}

Every entry of the last four columns is at the $99.3$rd percentile or above, every odd set falls below the $29$th, and
the two discriminants that showed nothing under the weight show a clean signal without it. So the apparent rule at
$p=3$ is an artefact of $\lambda$; nothing arithmetic is hidden there. Note also that the fit improves monotonically
along the three columns, for every one of the six discriminants: the more of the arithmetic weight one removes, the
more of the variance the even spectrum explains. That is itself evidence that the spectral term is real and that the
weight is what obscures it.

The last column is the cleanest object of all. Dropping both the weight and the squarefree condition leaves
\begin{equation}\label{eq:model}
\mathcal G(Y)=\sum_{h\le Y}(Y-h)\bigl(\sigma^{*}_{-1}(h^2-D)-\mathbb E\sigma^{*}_{-1}\bigr),\qquad
\sigma^{*}_{-1}(n)=\sum_{\substack{d\mid n\\ (d,2D)=1}}\frac1d,
\end{equation}
a Riesz mean of the divisor-type sums of Hooley and Gafurov. The condition $(d,2D)=1$ is retained; it is what makes
the levels $e\mid\operatorname{rad}(2D)$ appear (Step~2 of \S\ref{sec:formula} and Remark~\ref{rem:coprime}). Lemmas~\ref{lem:sawtooth} and~\ref{lem:weyl} never used
$\lambda$ or squarefreeness, so they apply verbatim, and the Dirichlet series attached to~\eqref{eq:model} is
\[
Z_k(s)=\sum_{(d,2D)=1}\mathcal W_k(D;d)\,d^{-s},
\]
a Sali\'e zeta function carrying a coprimality condition and nothing else. This is the object of
\S\ref{sec:formula}. The piece the off-diagonal of part~III actually requires is~\eqref{eq:model} perturbed twice:
by the squarefree condition, through M\"obius over $\ell^2$, and by $\lambda=1*\kappa$ with $\kappa(p)=4/(p-4)=O(1/p)$.

\paragraph{The phases: a test with no free parameter.} The sharpest check available to us is not the size of the
oscillation but its phase. In~\eqref{eq:orbit} the length enters the test function only through the combination
$t/\sqrt{|D|}$, so the spectral term carries the factor $(\sqrt{|D|}/2t)^{it_j}$ of~\eqref{eq:asymp}, that is
$(\sqrt{|D|}/2)^{it_j}$ in its dependence on $D$; and by Step~3 the amplitude is a real number, a signed combination of
periods of the real-valued form $u_j$, whose \emph{sign} therefore shifts the phase by $\pi$. (The sign used below is
that of the Katok--Sarnak period at $D$, as in the first version of this paper; the corrected Step~3 makes the
coefficient a combination over $e\mid\operatorname{rad}(2D)$ in which the level-one term is expected to dominate, and
the agreement found here is evidence that it does.) Writing
the fitted oscillation as $A_D\cos(t_1\log Y-\varphi_D)$, the mechanism predicts, once one discriminant is fixed as a
reference and with nothing else at our disposal,
\begin{equation}\label{eq:phasepred}
\varphi_D-\varphi_{D_0}=\frac{t_1}2\log\frac{|D|}{|D_0|}-\pi\cdot\bigl[\operatorname{sign}\mathrm{Per}_D\ne\operatorname{sign}\mathrm{Per}_{D_0}\bigr]\pmod{2\pi}.
\end{equation}
The signs of the periods are computed from the Fourier coefficients of the first even form, independently of the
arithmetic data. Taking $D_0=-4$ (\texttt{scripts/phase-test.py}):

\begin{table}[htbp]
\centering\small
\caption{The phase of the $t_1=13.7798$ oscillation of the model object, predicted by~\eqref{eq:phasepred} with no
free parameter, against the phase fitted to the data ($Y\le10^7$, $u=1$). The reference $D_0=-4$ is fitted, the other
five are predictions.}\label{tab:phases}
\begin{tabular}{lrrrr}
\toprule
$D$ & $\operatorname{sign}\mathrm{Per}_D$ & predicted $\varphi_D$ & observed $\varphi_D$ & difference\\
\midrule
$-3$  & $+$ & $+1.158$ & $+1.050$ & $+0.108$\\
$-4$  & $+$ & $+3.140$ & $+3.140$ & (reference)\\
$-7$  & $+$ & $+0.712$ & $+0.500$ & $+0.212$\\
$-8$  & $-$ & $-1.509$ & $-1.430$ & $-0.079$\\
$-11$ & $-$ & $+0.685$ & $+0.530$ & $+0.155$\\
$-19$ & $-$ & $-1.833$ & $-1.870$ & $+0.037$\\
\bottomrule
\end{tabular}
\end{table}

Every one of the five predictions is correct to within $0.212$ radians, that is to $3.4\%$ of a full period. Five
independent phases falling that close to prediction by chance has probability about $1.4\cdot10^{-6}$. The test uses
both halves of the mechanism at once: the factor $(2/\sqrt{|D|})^{it_1}$, which is what makes the predicted phases
differ from one another at all, and the sign of the Katok--Sarnak period, which is what puts $D=-8,-11,-19$ on the
opposite side. We regard this, rather than any of the variance fits, as the evidence that the mechanism of
\S\ref{sec:formula} is the right one.

For the amplitudes the agreement is weaker but not absent. Dividing the fitted amplitude by
$|\mathrm{Per}_D(u_1)|/|\Gamma_{z_D}|$ gives, in units of $10^7$, the values $2.68$, $2.22$, $0.39$, $1.90$, $2.71$,
$2.71$ for $D=-3,-4,-7,-8,-11,-19$: five of the six agree to within $40\%$ with no further factor, and the outlier is
$D=-7$, which is also the case whose fitted amplitude is least stable, moving by a factor $1.7$ between the two halves
of the range.

\subsection{The same test on Hooley's sum}\label{sec:hooley}

Everything above concerns a divisor sum smoothed by a Riesz mean. The identity~\eqref{eq:poincare} predicts the same
law for a barer object, namely the partial sums of the Weyl sums themselves,
\[
T_k(X)=\sum_{\substack{d\le X\\ (d,2D)=1}}\ \sum_{b\bmod d,\ b^2\equiv D}\e\Bigl(\frac{kb}d\Bigr),
\]
the sum Hooley bounded by $O(X^{3/4+\varepsilon})$ in \cite{Hooley1963}. Since the poles of its Dirichlet series are
those of~\eqref{eq:poincare}, at $s=\tfrac12\pm it_j$, the prediction is $T_k(X)=\sqrt X\sum_jc_jX^{it_j}+\cdots$, with
no smoothing anywhere. We computed $T_1(X)$ for all $X\le10^6$ (\texttt{scripts/weyl-partial.ts}; roots of
$b^2\equiv D$ by Tonelli--Shanks, Hensel and the Chinese remainder theorem).

$T_1(X)/\sqrt X$ is indeed bounded, staying below $0.56$ in absolute value with standard deviation $0.14$ for
$D=-4$: square-root cancellation, well beyond the proved exponent $\tfrac34$. Its spectrum is the even one, for every
discriminant tested, at the $96$th percentile or above, while the frequencies $2t_j$ are rejected at the $12.7$th.
And the phase test of~\eqref{eq:phasepred} applies verbatim, with $(\sqrt{|D|}/2)^{it_1}$ entering through the
factor $(\sqrt{|D|}/2t)^{it_1}$ of~\eqref{eq:asymp}:

\begin{table}[htbp]
\centering\small
\caption{The same parameter-free phase prediction~\eqref{eq:phasepred}, applied to the partial sums of Hooley's Weyl
sums $T_1(X)$, $X\le10^6$. Reference $D_0=-4$; the other five are predictions.}\label{tab:weylphase}
\begin{tabular}{lrrrrr}
\toprule
$D$ & $R^2$ even & percentile & predicted $\varphi_D$ & observed $\varphi_D$ & difference\\
\midrule
$-3$  & $0.173$ & $100$  & $-2.992$ & $-2.980$ & $-0.012$\\
$-4$  & $0.412$ & $100$  & $-1.010$ & $-1.010$ & (reference)\\
$-7$  & $0.159$ & $100$  & $+2.846$ & $+3.110$ & $-0.264$\\
$-8$  & $0.252$ & $98.3$ & $+0.624$ & $+0.520$ & $+0.104$\\
$-11$ & $0.170$ & $100$  & $+2.818$ & $+2.810$ & $+0.008$\\
$-19$ & $0.206$ & $96.0$ & $+0.301$ & $+0.310$ & $-0.009$\\
\bottomrule
\end{tabular}
\end{table}

Three of the five predictions are correct to within $0.012$ radians, two parts in a thousand of a full period, and
the median discrepancy over the five is $0.012$. We know of no way to obtain agreement of this kind other than by the
mechanism of \S\ref{sec:formula} being correct.

\paragraph{The limits of the present computations, and what is planned.} Every number reported above comes from a
sieve of $Q_u(h)$ for $h\le10^7$, which runs in seconds to minutes on a laptop; the scripts are in the repository and
each states its cost in its header. Three of the questions raised here are limited by that range rather than by
anything conceptual, and are settled by longer runs of exactly the same code:
\begin{enumerate}
\item the amplitude law of Table~\ref{tab:periods}, whose comparison is currently blunted by an uncertainty of up to a
factor $2$ in the fitted amplitudes; the uncertainty falls like $Y^{-1/2}$ in the length of the $\log Y$ range, so
$Y=10^9$ would roughly halve it and $Y=10^{12}$ would settle the discrepancy at $D=-7$;
\item the resolution in the frequency, which is $2\pi/\log(Y_{\max}/Y_{\min})\approx0.68$ at present and therefore
cannot separate the lines $19.42$ and $19.48$, the first even and the first odd parameter of that size; this is the
single measurement that would test the parity rule against an adjacent pair rather than against a whole set;
\item the selection rule at $p=3$, where the sample is fourteen polynomials and should be a few hundred.
\end{enumerate}
None of the three needs a new method, only more arithmetic, and we intend to carry the computation out on
appropriate hardware and report it separately. We state them here so that a reader can see exactly which of our
assertions rest on $10^7$ and which do not: the parity rule of Table~\ref{tab:fits} and the threshold of
Table~\ref{tab:periods} are stable across the halves of the present range and do not, the amplitude law does.

\section{Uniformity in $u$, and what the Ces\`aro conjecture would need}\label{sec:uniformity}

Theorem~\ref{thm:main} is for one piece, $u=1$. Part~III reduces the Ces\`aro form of Conjecture~1 of part~I to the
windows of \emph{all} the pieces with $u\le H^{2/3+\varepsilon}$, so what is wanted is Theorem~\ref{thm:main} uniformly
in $u$, with an error that survives summation against the weight $w(u)$ over that range. This section records what
that requires, and why the proof above does not supply it.

For $u>1$ the roots are those of $u^2x^2\equiv D$, and by Lemma~\ref{lem:bridge} the dilation moves the frequency
from $k$ to $k\bar u$, a frequency depending on the modulus. Geometrically, the pairs $(d,x)$ with
$u^2x^2\equiv D\ (d)$ are again Heegner points, now of discriminant $4u^2D$: the associated form is
$[d,2u^2x,\ast]$ and the point is
\[
z=\frac{-u^2x+u\sqrt D}{d},\qquad \Im z=\frac{u\sqrt{|D|}}d,\qquad \Re z=-\frac{u^2x}d .
\]
Here is the difficulty, in the sharpest form we can state it. Since $x$ runs modulo $d$, the coordinate $\Re z$ runs
modulo $u^2$, not modulo $1$: the seed is invariant under $z\mapsto z+u^2$ but not under $z\mapsto z+1$. The relevant
Poincar\'e series is therefore taken over $\Gamma_\infty^{(u^2)}\backslash\Gamma$, where $\Gamma_\infty^{(u^2)}$ is
generated by $z\mapsto z+u^2$, and in the pairing with $u_j$ only the Fourier modes with $k\equiv0\ (u^2)$ survive.
Consequently the geometric factor is \emph{not} the full Katok--Sarnak period over the Heegner points of
discriminant $4u^2D$, but a sum over a restricted sub-family; the relation of Shimura, which would express the full
period at $Df^2$ through Hecke eigenvalues at $f$ and so cost only $u^{7/64+\varepsilon}$, does not apply to it. This is
exactly the dilation obstruction of part~III, seen geometrically.

What does survive is the analysis of \S\ref{sec:proof}. Carrying it through with the modes $k=nu^2$ gives
$\alpha=2\pi n$ and $\beta=2\pi nut/\sqrt{|D|}$, so the expansion parameter of Lemma~\ref{lem:mellin} becomes
\begin{equation}\label{eq:ugeneral}
\frac\alpha\beta=\frac{\sqrt{|D|}}{ut},
\end{equation}
which reduces to the $u=1$ case and predicts, through the factor $(\alpha/2\beta)^{it_j}$, that the phase of the
$t_j$-oscillation shifts by exactly $-t_j\log u$ relative to $u=1$. That is a prediction with no free parameter, and
it holds: for $u=2$ and the five discriminants $D=-4,8,-8,-7,12$ the differences between predicted and observed
phases are $0.002$, $-0.118$, $0.305$, $3.005$ and $0.302$ radians, so four of the five agree to within $0.31$ and
the fifth, which has the least stable fit of the six, differs by $\pi$, i.e. by a sign in the geometric factor.\footnote{%
Computed from the model object at $Y\le10^7$ with the mean value evaluated to $22$ digits
(\texttt{scripts/ef-general.py}); the median discrepancy over the five is $0.302$ radians.}

So the structure extends to every $u$, and what is missing is precisely one thing: the size, in $u$, of the
restricted geometric factor. Summing $\lambda\omega(u)$ times an error $(H/u)^{1/2}u^{A}$ over $u\le H^{2/3}$ is
$O(H^{1-\delta})$ exactly when $A<\tfrac14$, so a bound $u^{1/4-\varepsilon}$ for that factor would settle the Ces\`aro
conjecture for quadratics. This is the same wall part~III meets from the other side, in the guise of the exponent
condition $\theta+6B<1$: there, no Weyl-sum bound is uniform enough in the frequency; here, no bound is known for the
restricted period. We regard the two as one open problem seen from two directions, but the present formulation is
sharper, because the object to be bounded is now a specific sum of a Maass form over a specific set of Heegner
points, rather than an error term.

What the present paper does say about it is that the object on the far side is now identified. The windows of
part~III are not an amorphous error term: piece $u$ carries the spectrum of level $4u^2$, with amplitudes the
periods at discriminant $4u^2D$, and the Ces\`aro conjecture is the statement that these spectra, summed over $u$
with the weight $w(u)/u$, exhibit enough cancellation. That is a sharper question than the one part~III could pose,
and it is the one we would attack next.

\section*{AI-assisted research: disclosure}
This work was carried out with heavy use of large language models, the Anthropic models Claude Fable~5.1 and
Claude Opus~5, used interactively by the author in September 2026 as research assistants. The author posed the questions, set the
direction and the scope, and made every decision about what to pursue, claim and submit. The model was used to draft
mathematical arguments, to write and run the software behind every number in the paper, to search the literature,
and to draft the text; all of this was done in dialogue with the author, who verified the code and the numerical
results, reviewed the proofs, and takes full responsibility for the content. The model is not an author. The
complete record of the work is in the public repository \url{https://github.com/gewure/ulamnd} (directory \texttt{research/}).

\bibliographystyle{plain}
\bibliography{refs}
\end{document}
